\section{Machine verification in Lean 4}\label{sec:lean}

The theorems of this paper are statements of algebraic geometry and are not
formalised. What is formalised is the finite arithmetic on which many of them
turn, which the kernel can settle by computation; the rest of their
computational content is checked in exact arithmetic by the computations of
\Cref{sec:verification}. The certificate is a single Lean~4 file that
accompanies the paper. It compiles against a bare Lean~4 toolchain, contains
no \texttt{sorry}, and depends on no library, Mathlib included. The axiom
report at the end of the file lists, for each of its four hundred and nine theorems,
the axioms on which it depends. One hundred and twenty-eight depend on none, and the
remaining two hundred and eighty-one on propositional extensionality alone, which enters when
\texttt{decide} reduces an equality of Booleans or when a bound is proved from
the lemmas on natural numbers in the core library; no theorem uses anything
else. The rank-zero obstruction of \Cref{thm:rankzero} is not among the
four hundred and nine: it is proved by hand from \Cref{lem:annihilator}, whose count
the certificate checks in its model of Hochschild cohomology, and its
application to the explicit object is verified in exact arithmetic modulo a
prime by item (XIV) of \Cref{sec:verification}.

\subsection{Contents of the certificate}

\begingroup
\small
\renewcommand{\arraystretch}{1.06}
\begin{longtable}{@{}>{\raggedright\arraybackslash}p{0.31\linewidth}>{\raggedright\arraybackslash}p{0.62\linewidth}@{}}
\caption{The four hundred and nine machine-checked theorems.}\label{tab:lean}\\
\toprule
\textbf{theorem} & \textbf{statement} \\
\midrule
\endfirsthead
\multicolumn{2}{@{}l}{\footnotesize\itshape Table \thetable, continued.}\\[2pt]
\toprule
\textbf{theorem} & \textbf{statement} \\
\midrule
\endhead
\midrule
\multicolumn{2}{r@{}}{\footnotesize\itshape continued on the next page}\\
\endfoot
\bottomrule
\endlastfoot
\texttt{qnorm\_mul\_check} & the norm on $\ZZ[\sqrt{-d}]$ is multiplicative on the data used \\
\texttt{qmul\_conj} & $\tau\bbar{\tau}$ equals the norm, as an element of $\ZZ[\sqrt{-d}]$ \\
\texttt{characters\_differ} & $\tau^{2n}\ne\Nm(\tau)^{n}$ at $\tau=2+\sqrt{-d}$, for nine fields and $n\le8$ \\
\texttt{characters\_differ'} & the same at five further witnesses \\
\texttt{bad\_witness} & the witness $\tau=1+\sqrt{-d}$ fails, at $(d,n)=(1,4)$ and $(3,3)$ \\
\texttt{bad\_witness\_exceptions} & and fails exactly when $d=1,\,4\mid n$ or $d=3,\,3\mid n$ \\
\texttt{ratio\_not\_root\_of\_unity} & $\tau^{n}\ne\bbar{\tau}^{\,n}$ at $\tau=2+\sqrt{-d}$ \\
\texttt{sign\_uniform\_and\_closed} & the sign of \Cref{thm:fmtheta} is
  uniform over index sets and equals $(-1)^{g(g+1)/2+k}$, for $g\le5$ \\
\texttt{subsets\_count} & the subset enumeration used there has $\binom{g}{k}$ members \\
\texttt{semiregularity\_target} & $\sum_{q}\binom{2n}{q}\binom{2n}{q+2}=\binom{4n}{2n-2}$ for $n\le7$ \\
\texttt{secant\_count\_}\allowbreak\texttt{below\_thresholds} &
  $2n^{2}+2n<2^{2n}<3^{2n}$ for $2\le n\le8$ \\
\texttt{thresholds\_not\_necessary} & $3<2^{2}$ and $5<3^{2}$, which show that neither threshold of \Cref{sec:numerology} is necessary \\
\texttt{grading\_positions\_differ} & $(2n,0)$, $(0,2n)$ and $(n,n)$ are distinct cells for $n\ge1$ \\
\texttt{weil\_delta\_parity} & $\sqrt{-d}^{\,k}$ is rational exactly when $k$ is even, for $k\le20$ \\
\texttt{weil\_delta\_closed\_even} & $\sqrt{-d}^{\,2m}=(-d)^{m}$ \\
\texttt{weil\_delta\_closed\_odd} & $\sqrt{-d}^{\,2m+1}=(-d)^{m}\sqrt{-d}$ \\
\texttt{weil\_term\_counts} & each of $\omega_{1},\omega_{2}$ has $2^{2n-1}$ monomials \\
\texttt{supports\_disjoint} & no monomial of the Weil line occurs in a power of $\eta$ \\
\texttt{hodge\_dim\_count} & $\dim\Hdg^{k}=3$ for $k=n$ and $1$ otherwise, for $n\le6$ \\
\texttt{weil\_line\_hodge\_type} & the Weil line is of type $(n,n)$ exactly at balanced signature \\
\texttt{clifford\_dimensions} & $\dim C=2^{b}$, $\dim C^{+}=2^{b-1}$, $\dim\mathrm{KS}=2^{b-2}$ \\
\texttt{k3\_type\_only\_at\_one} & $h^{2,0}=1$ on the pieces of $H^{2}$ forces $n=1$ \\
\texttt{moduli\_codimension} & the Weil family has codimension $n(n+1)$ in $\cA_{2n}$ \\
\texttt{discriminant\_is\_a\_norm} & $d^{n}p^{2}$ is a norm from $K$, with an explicit witness \\
\texttt{split\_codimension} & the split locus has codimension $n(n-1)/2$, zero only at $n=1$ \\
\texttt{quat\_class\_odd} & for $n$ odd the quaternionic class is $b$ times a square \\
\texttt{quat\_class\_even} & for $n$ even it is a square \\
\texttt{quat\_product\_class} & $d^{\,n-1}$ is a norm for $n$ even, with a witness \\
\texttt{quat\_locus\_codimension} & the quaternionic locus has dimension $n(n+1)/2$, codimension $n(n-1)/2$ \\
\texttt{quat\_locus\_}\allowbreak\texttt{divisor\_at\_two} & at $n=2$ that locus is a divisor: the three numbers are $4$, $3$, $1$ \\
\texttt{irrational\_}\allowbreak\texttt{witness\_exists} & some $m\le7$ makes $(m+\sqrt{-d})^{2n}$ irrational, for $n\le8$ \\
\texttt{weil\_line\_spanned} & the determinant of $w,\tau w$ is $b(w_{1}^{2}+dw_{2}^{2})$, hence nonzero \\
\texttt{norm\_form\_positive} & $a^{2}+db^{2}>0$ for every nonzero $(a,b)$, for ten fields \\
\texttt{norm\_sum\_}\allowbreak\texttt{forces\_zero} & a positive combination of norms vanishes only when every term does \\
\texttt{semiregularity\_}\allowbreak\texttt{source\_target} & at $n=2$ the source of $\sigma$ overtakes its target at $s=5$ \\
\texttt{tensor\_rank\_gap} & $\binom{2n}{2}\ge6>1$ for $n\ge2$, the separation of \Cref{lem:tensorrank} \\
\texttt{secant\_witness\_}\allowbreak\texttt{at\_n\_four} & the line bundle witness of \Cref{thm:secantexists} reproduces $6(u_{t}+v_{t})$ in every degree \\
\texttt{secant\_witness\_rank} & that witness has rank $6=Ma$, which is positive, so \Cref{lem:posrank} applies \\
\texttt{vandermonde\_}\allowbreak\texttt{nodes\_distinct} & the Vandermonde product on the nodes $0,\dots,n$ is nonzero for $n\le8$ \\
\texttt{level\_sums\_forced} & the level sums of \Cref{thm:pterange} satisfy $\sum_{j}M_{j}N_{j}^{r}=0$, $r=1,2,3$, at six level sets \\
\texttt{level\_sums\_totals} & their totals are $1$, $3$ and $-2$, so none of those objects has rank zero \\
\texttt{level\_matrix\_}\allowbreak\texttt{nonsingular} & with at most $n$ distinct norms the matrix $(N_{j}^{r})$ is nonsingular \\
\texttt{gauss\_norm\_counts} & the number of Gaussian integers of norm $1,2,3,4,5,9,45$ \\
\texttt{norm\_supply\_}\allowbreak\texttt{blocks\_small\_case} & at the nodes $1,2,3,4$ the forced $|M_{3}|=4$ while $\ZZ[i]$ has no element of norm $3$ \\
\texttt{pencil\_volume} & $\int\eta^{4}=4!\,(2d)^{4}=384d^{4}$ on $\cO^{2}$, as a polynomial identity in $d$, \Cref{prop:pencil} \\
\texttt{pencil\_square} & $C_{k}^{2}=(1+k^{2})I$ for the matrix of the pencil, so $\phi_{k}^{2}=(1+k^{2})/(4d^{2})$ \\
\texttt{pencil\_minimum} & $\tfrac{8}{24}\beta(k)\int\eta^{4}=32d^{2}(1+k^{2})$, the minimum $\mu_{k}$ \\
\texttt{pencil\_degree\_bound} & $\mu_{k}\int\eta^{4}/2=6144d^{6}(1+k^{2})$ and $6144=32^{2}\cdot6$, the degree bound of \Cref{prop:pencil} \\
\texttt{split\_obstruction\_}\allowbreak\texttt{rank} & $n^{2}-n(n+1)/2=n(n-1)/2$ and $n(n-1)=2\cdot n(n-1)/2$ for $n\le30$: the rank of the obstruction map of \Cref{prop:splitobstruction}(ii) and the even side of \Cref{cor:pteobstructed} \\
\texttt{split\_obstruction\_}\allowbreak\texttt{count} & the bookkeeping of \Cref{prop:splitobstruction}(iv) at $n=3$: kernel $15=9+6$, image $3$, direct sum $12$, complement $3$ \\
\texttt{evaluation\_not\_}\allowbreak\texttt{surjective} & $s\binom{2n}{2}>2\binom{2n}{2}+4n^{2}$ for $5\le s\le40$ and $2\le n\le30$, so the evaluation map of \Cref{prop:evaluation}(iii) is not surjective; at $n=3$, $s=6$ the dimensions are $66$ and $90$ \\
\texttt{markman\_candidate\_}\allowbreak\texttt{identities} & for odd $d\le201$ and $N=(d+9)/2$: the point count $12N(N-5)=(d+9)(3d-3)$, the Euler characteristic $(d+9)(27-d)$, the two identities that place $\ch(\cF)$ at $(1,3)$, $\Nm(3+\sqrt{-d})=2N$, and the nonvanishing of the ranks, \Cref{prop:markman8}\\
\texttt{secant\_kernel\_}\allowbreak\texttt{arithmetic} & the matrix $\bigl(\begin{smallmatrix}a&b\\ b&-ad\end{smallmatrix}\bigr)$ has determinant $-(a^{2}d+b^{2})\ne0$ for $(a,b)\ne(0,0)$, $a,b<13$, $d\le20$; the coefficients $C_{k}=k!\,c_{k}$ of $au_{t}+bv_{t}$ satisfy $C_{k+1}=-d\,C_{k-1}$; and $n(n+1)/2+n(n-1)/2=n^{2}$ for $n<60$, \Cref{thm:factor}\\
\texttt{candidate\_unnormalised\_}\allowbreak\texttt{points} & for odd $d\le201$: $(d+9)(3d-3)-(d+9)(2d-1)=(d+9)(d-2)\ge12$, so an unnormalised double point exists; normalising all of them gives $\chi=50N$, and $50N\ne72N-4N^{2}$ for every integer $N>0$, \Cref{thm:candidatefails}, \Cref{rem:normalisedpoints}\\
\texttt{mumford\_rigidity\_}\allowbreak\texttt{sos} & a positive multiple of each of the seventeen forms $P_{1},P_{2},Q_{1},\dots,Q_{4},R,S_{1},\dots,S_{4},Z_{1},\dots,Z_{6}$ is a sum of squares of integral linear forms with positive coefficients, checked on $\{0,1,2\}^{4}$, which suffices for polynomials of degree at most two in each variable, \Cref{thm:mumfordrigid} \\
\texttt{smooth\_invariants} & Noether's formula and the self-intersection formula for a smooth support, $K^{2}+e=12\chi$ and $K^{2}-e=6(b^{2}+d)^{2}$ \\
\texttt{smooth\_bmy\_defect} & the Bogomolov--Miyaoka--Yau defect of a smooth support is $24(b^{4}-d^{2})$ \\
\texttt{smooth\_bmy\_}\allowbreak\texttt{is\_d\_le\_bsq} & so the inequality is exactly $d\le b^{2}$, over a box in $b$ and $d$ \\
\texttt{smooth\_index\_vacuous} & the Hodge index bound $9(b^{2}+d)\ge8b^{2}$ holds for every $d\ge1$ \\
\texttt{smooth\_four\_}\allowbreak\texttt{discriminants} & the squarefree odd $d\le9$ are exactly $1,3,5,7$, which is \Cref{cor:fourdiscriminants} \\
\texttt{smooth\_table} & the invariants of \Cref{tab:smoothsupport} at $b=3$ \\
\texttt{burch\_weight\_}\allowbreak\texttt{identities} & the weight $x^{2}+d$ annihilates three moments, $c_{k+2}+dc_{k}=0$, which is \Cref{thm:noci} \\
\texttt{burch\_rank\_one\_}\allowbreak\texttt{discriminant} & nine times the discriminant of the rank one quadratic is $-2b^{2}-18d$, always negative \\
\texttt{burch\_rank\_four\_}\allowbreak\texttt{witness} & the rank four candidate at $d=23$ solves the moment system and is removed only by injectivity \\
\texttt{closure\_omits\_}\allowbreak\texttt{conjecture} & the Hodge conjecture is not in the forward closure of what this paper proves and quotes \\
\texttt{frontier\_}\allowbreak\texttt{suffices} & it is in the closure once (F3) alone is adjoined, which is \Cref{prop:f3ishc}; once $\{\mathrm{L},\mathrm{M}\}$, $\{\mathrm{L},\mathrm{F3}'\}$, $\{\mathrm{V},\mathrm{F3}'\}$ or $\{\mathrm{P2}_{\mathrm{split}},\mathrm{F2},\mathrm{F3}'\}$ is, with $\mathrm{P2}_{\mathrm{split}}$ the propagation statement for the split families of the CM fields of degree at least four; and, not minimally, once $\{\mathrm{P2}_{\mathrm{split}},\mathrm{F2},\mathrm{F3}\}$, $\{\mathrm{L},\mathrm{F3}\}$ or $\{\mathrm{V},\mathrm{F3}\}$ is \\
\texttt{frontier\_}\allowbreak\texttt{minimal} & in each of the five minimal sets of \Cref{thm:closure}(ii) every element is necessary: dropping any one leaves the conjecture underivable \\
\texttt{frontier\_}\allowbreak\texttt{smallest} & of the sixteen open statements exactly one, (F3), suffices alone, and of the one hundred and twenty pairs exactly those containing (F3) and the pairs $\{\mathrm{L},\mathrm{M}\}$, $\{\mathrm{L},\mathrm{F3}'\}$, $\{\mathrm{V},\mathrm{F3}'\}$ do \\
\texttt{variational\_gives\_}\allowbreak\texttt{abelian} & (V) alone puts the conjecture for abelian varieties in the closure, and not $\mathbf{HC}$, \Cref{thm:vhcab} \\
\texttt{lefschetz\_gives\_}\allowbreak\texttt{abelian} & (L) gives (V), hence the conjecture for abelian varieties, and not $\mathbf{HC}$, \Cref{prop:lefvhc} \\
\texttt{secant\_route\_}\allowbreak\texttt{descends} & granting both open demands of the secant route reaches $\delta=\delta_{0}$ and, through \Cref{prop:descent}, every imaginary quadratic family, but neither the CM fields of higher degree nor the conjecture \\
\texttt{propagation\_}\allowbreak\texttt{closes} & (P2) for the split imaginary quadratic families puts the whole imaginary quadratic case in the closure, and (P2) for the split families of the fields of degree at least four puts the Weil classes of every CM field there \\
\texttt{p2\_minimal\_}\allowbreak\texttt{count} & $\binom{4n}{2}-4n^{2}=2\binom{2n}{2}=2n(2n-1)$ for every $n\le60$, the dimension at which the pure form of the criterion would be met, \Cref{thm:p2numerical}, and which no complex attains, \Cref{thm:p2false} \\
\texttt{mumford\_}\allowbreak\texttt{invariants} & the invariants of $\mathfrak{sl}_{2}^{3}$ in $\bigwedge^{q}(V_{1}\otimes V_{2}\otimes V_{3})$ have dimensions $1,0,1,0,1,0,1,0,1$, counted from weight multiplicities, \Cref{cor:mumfordlefschetz} \\
\texttt{two\_branch\_}\allowbreak\texttt{annihilator} & for $n=2,3$ the monomials of type $(p,p)$ killed by $Q\wedge HH^{1}$, $P\wedge HH^{1}$ and $HH^{1}\wedge HH^{1}$ are $\alpha_{-}$, $\alpha_{+}$ and none, \Cref{thm:nosum} \\
\texttt{weil\_monomials\_}\allowbreak\texttt{separated} & neither Weil monomial lies in $\bigwedge^{*}(V_{-}^{1,0}\oplus V_{+}^{0,1})$ or $\bigwedge^{*}(V_{+}^{1,0}\oplus V_{-}^{0,1})$, for $n\le12$, which is the last step of \Cref{thm:p2support}(iii) \\
\texttt{mumford\_ext\_}\allowbreak\texttt{profile} & the bounds $1,16,119,328,560,328,119,16,1$ are palindromic, with alternating sum $112$, sum $1488$ and $2+238+560=800$, \Cref{thm:mumfordprofile} \\
\texttt{lefschetz\_}\allowbreak\texttt{counts} & the $\Sp_{8}$-invariants of $\bigwedge^{*}(V\oplus V)$ number $1,3,6,10,15,10,6,3,1$ by the decomposition of each $\bigwedge^{i}V$, the Hodge classes that are not Lefschetz $0,0,2,6,13,6,2,0,0$, the Lefschetz classes at a CM point are counted by $(1+4y+y^{2})^{4}$, and Weyl's formula gives $308$ for $2\varpi_{2}$, \Cref{thm:lefschetzclosure} \\
\texttt{criterion\_}\allowbreak\texttt{endomorphisms\_n2} & for all natural numbers, $2e_{0}+12=\chi+2e_{1}$ with $e_{1}\ge8$ and $\chi\ge1$ gives $e_{0}\ge3$: the Euler characteristic of an object meeting the criterion at $n=2$, \Cref{thm:p2endo}(iv) \\
\texttt{criterion\_}\allowbreak\texttt{endomorphisms\_n3} & for all natural numbers, $2e_{0}+60=\chi+2e_{1}+e_{3}$ with $e_{1}\ge12$, $e_{3}\ge40$ and $\chi\ge1$ gives $e_{0}\ge3$: the same at $n=3$, \Cref{thm:p2endo}(iv) \\
\texttt{quartic\_}\allowbreak\texttt{squares\_}\allowbreak\texttt{mod4\_}\allowbreak\texttt{phi} & the squares modulo $4$ in $\ZZ[\varphi]$, $\varphi^{2}=\varphi+1$, are $0$, $1$, $1+\varphi$, $2+3\varphi$, \Cref{prop:quarticother} \\
\texttt{quartic\_}\allowbreak\texttt{no\_}\allowbreak\texttt{square\_}\allowbreak\texttt{phi} & neither $-1$ nor twice a unit is a square modulo $4$ in $\ZZ[\varphi]$, \Cref{prop:quarticother} \\
\texttt{quartic\_}\allowbreak\texttt{squares\_}\allowbreak\texttt{mod4\_}\allowbreak\texttt{sqrt2} & the squares modulo $4$ in $\ZZ[\sqrt2]$ are $0$, $1$, $2$, $3+2\sqrt2$, \Cref{prop:quarticother} \\
\texttt{quartic\_}\allowbreak\texttt{no\_}\allowbreak\texttt{square\_}\allowbreak\texttt{sqrt2} & $3$, $1+3\sqrt2$ and $1+\sqrt2$ are not squares modulo $4$ in $\ZZ[\sqrt2]$, \Cref{prop:quarticother} \\
\texttt{quartic\_}\allowbreak\texttt{rank\_}\allowbreak\texttt{two\_}\allowbreak\texttt{congruences} & $m\equiv9$ has no solution $1\le m\le8$ modulo $13$ or $17$, and modulo $5$ only $m=4$, \Cref{prop:quarticother} \\
\texttt{quartic\_}\allowbreak\texttt{euler\_}\allowbreak\texttt{minimal} & the minimal profiles $1,8,r,8,1$ with $r=18,20,12$ have Euler characteristic $4,6,-2$, \Cref{thm:quarticobstruction} \\
\texttt{quartic\_}\allowbreak\texttt{squares\_}\allowbreak\texttt{mod4\_}\allowbreak\texttt{mod8} & a square is $0$ or $1$ modulo $4$, an odd square is $1$ modulo $8$, and twice a square is $0$ or $2$ modulo $8$, \Cref{prop:quarticother} \\
\texttt{quartic\_}\allowbreak\texttt{rank\_}\allowbreak\texttt{two\_}\allowbreak\texttt{small\_}\allowbreak\texttt{odd\_}\allowbreak\texttt{parts} & for $D=3,7$ no $m\in[1,8]$ is $9$ modulo $D$ and $0$ or $1$ modulo $4$; for $D=6,10,14$ none is $9$ modulo $D/2$ and $1$ modulo $4$; for $D=2$ only $m=1,5$ are $1$ modulo $4$, \Cref{prop:quarticother} \\
\texttt{quartic\_}\allowbreak\texttt{rank\_}\allowbreak\texttt{two\_}\allowbreak\texttt{large\_}\allowbreak\texttt{odd\_}\allowbreak\texttt{part} & for every modulus $n\ge10$ no $m<9$ is $9$ modulo $n$, \Cref{prop:quarticother} \\
\texttt{quartic\_}\allowbreak\texttt{sqrt2\_}\allowbreak\texttt{units\_}\allowbreak\texttt{mod4} & $3$ and $1+2\sqrt2$ are not squares modulo $4$ in $\ZZ[\sqrt2]$, $(1+\sqrt2)^{2}\equiv3+2\sqrt2$ and $(3+2\sqrt2)^{2}\equiv1$ modulo $4$, and $5\not\equiv\pm1$ modulo $8$, \Cref{prop:quarticother} \\
\texttt{sextic\_}\allowbreak\texttt{thresholds} & the thresholds $r^{3}-2r^{2}+22$ of the six computed sextic profiles are $2$, $-6$, $10$, $26$, $8$, $12$, a minimal object with the generic profile has $2(1-12+54)-112=-26$, and $(-4)^{3}=-64$, \Cref{prop:sexticcount} \\
\texttt{orlov\_}\allowbreak\texttt{equality\_}\allowbreak\texttt{count} & $n(n-1)+(2n)^{2}+n(n-1)=6n^{2}-2n$ for $3\le n\le64$, and $1+16+1=18$, \Cref{prop:orlovequality} \\
\texttt{orlov\_}\allowbreak\texttt{n4\_}\allowbreak\texttt{euler} & the minimal profile at $n=4$ has $\chi=-2$, while every secant character has $\chi=8d(a^{2}d+b^{2})\ge8$, \Cref{prop:orlovfour} \\
\texttt{cm\_tetrahedra} & the zero-sum four-sets of weights at a CM point are the six unions of two opposite pairs and the two tetrahedra $T_{+}$, $T_{-}=-T_{+}$, with even and odd numbers of entries $-1$, each meeting every opposite pair once, \Cref{prop:cmsource}(i) \\
\texttt{cm\_square\_}\allowbreak\texttt{monomials} & on the square, $132$ monomials of weight zero, $100$ with an opposite pair and $32$ tetrahedron monomials, $16$ for each tetrahedron, distributed $1,4,6,4,1$ by the number of factors from the first copy, \Cref{prop:cmsource}(ii) \\
\texttt{cm\_tetrahedron\_}\allowbreak\texttt{pairs} & the six two-element subsets of $T_{+}$ differ in exactly two coordinates, a subset and its complement have the same difference set, and each of $\{1,2\}$, $\{1,3\}$, $\{2,3\}$ occurs for exactly two subsets, \Cref{prop:cmsource}(ii) \\
\texttt{cm\_profile\_}\allowbreak\texttt{forces\_equal} & on $[-10,10]^{3}$, equal values of $|a_{2}-a_{3}|$, $|a_{1}-a_{3}|$, $|a_{1}-a_{2}|$ force $a_{1}=a_{2}=a_{3}$, and two of them agree exactly along $a_{1}=a_{2}$ or $a_{1}+a_{2}=2a_{3}$ and the two analogous pairs of relations, \Cref{prop:cmsource}(ii) and (iii) \\
\texttt{cm\_pairs\_}\allowbreak\texttt{three\_cycle} & a permutation of the four diagonals of the cube that fixes none of the three pairings of the diagonals is one of the eight cycles of length three, of order three with exactly one fixed diagonal, \Cref{lem:splitcm} \\
\texttt{cm\_abelian\_}\allowbreak\texttt{semiregular\_}\allowbreak\texttt{count} & for $n\le60$ the annihilator of the class of an abelian subvariety of half dimension has dimension $2\binom{2n}{2}-2\binom{n}{2}+3n^{2}=6n^{2}-n$, and its complement in $HT^{2}$ has dimension $\binom{2n}{2}$, that of $\Ext^{2}(\cO_{B},\cO_{B})$, \Cref{lem:absemireg} \\
\texttt{quartic\_rank\_}\allowbreak\texttt{values} & the tuples $(\mu,\rho_{1},\rho_{2},R_{1},R_{2})$ allowed at a quartic CM field give exactly fifteen values of $r(\gamma)$, the least $80$ only for $(0,0,0,8,8)$, never $100$, and nine when $\rho_{1}=\rho_{2}$, \Cref{thm:quarticrank}, \Cref{cor:quarticleast} \\
\texttt{delsarte\_loop\_}\allowbreak\texttt{sextic} & for the loop sextic $AB=15624\,I$ with $B$ the circulant $(3125,-625,125,-25,5,-1)$, whose rows sum to $2604$, and the Jacobian ring has dimensions $1$, $426$, $1751$, $426$, $1$ and $15625$ in all, \Cref{prop:delsarte} \\
\texttt{simplex\_klein\_}\allowbreak\texttt{quartic} & for the Klein quartic $x^{3}y+y^{3}z+z^{3}x$, $\det A=28$, the differences $R$ of its exponents have $\det R=7$ and $CR=7I$ for an integer matrix $C$ not divisible by $7$, so $e_{A}=7$, and the sum over the orbits gives $4$ times the Euler number $-4$, \Cref{cor:delsarteall} \\
\texttt{simplex\_loop\_}\allowbreak\texttt{sextic} & for the loop sextic $\det R=2604=2^{2}\cdot3\cdot7\cdot31$ and $CR=2604\,I$ with, for each of $2,3,7,31$, an entry of $C$ it does not divide, so $e_{A}=2604$; the sums over the orbits for the loop and the Fermat sextic are $6\cdot2610$, the Euler number of a smooth sextic fourfold, \Cref{lem:simplexchart}, \Cref{cor:delsarteall} \\
\texttt{quartic\_exceptional\_}\allowbreak\texttt{pair} & in $\bigwedge(V_{\sigma}\oplus V_{\bar\sigma})$ at a real place of a quartic CM field, $\theta^{4}=24\,\alpha_{\sigma}\alpha_{\bar\sigma}$ and $(\alpha_{\sigma}+\alpha_{\bar\sigma})^{2}=2\,\alpha_{\sigma}\alpha_{\bar\sigma}$; $\theta^{3}D=0$ and the four $\theta^{2}D$ are distinct monomials for $D$ in $V_{\sigma}^{0,1}\otimes V_{\bar\sigma}^{0,1}$; the two deformations that are not $F$-linear send $\alpha_{\sigma},\theta,\theta^{2}$ to $\theta b_{\sigma,0}b_{\sigma,1}$, $2b_{\bar\sigma,0}b_{\bar\sigma,1}$, $4\theta b_{\bar\sigma,0}b_{\bar\sigma,1}$ and symmetrically, so that a place is exceptional exactly when $16c^{2}=u_{\sigma}u_{\bar\sigma}$, \Cref{prop:quarticlocus} \\
\texttt{quartic\_locus\_}\allowbreak\texttt{counts} & the rational characters with both places exceptional have rank $94$ or $110$; the annihilator has dimension $8+4a+b\in\{8,9,10,12,13,16\}$; $\mathfrak{so}(4,3)$ has dimension $21$ and maximal compact subalgebra of dimension $9$, against $15$ and $7$ for $\mathfrak{su}(2,2)$, and the orbits have dimension $21-16=5$ and $15-11=4$, \Cref{prop:quarticlocus}, \Cref{cor:quarticlocus} \\
\texttt{quartic\_compensated\_}\allowbreak\texttt{bivectors} & in the exterior algebra on $H^{1}(X)=U_{1}\oplus U_{2}$, with $q=2$, the classes $x_{1}$, $x_{2}$ of \eqref{eq:quarticcompensated} annihilate the four classes spanning $S(0,2)$, and the six symmetric maps at the two places annihilate $\theta_{1}$ and $\theta_{2}$, \Cref{thm:quarticlocal}(i) \\
\texttt{quartic\_rank\_}\allowbreak\texttt{certificates} & a $20\times20$ minor of the contraction matrix of $4\beta'$ and a $12\times12$ minor of that of $\alpha_{0}$, for $f_{1}=2$, $f_{2}=\tfrac12$, are invertible modulo $1000003$, so the ranks are $20$ and $12$, \Cref{prop:markmanquarticclasses}(ii) \\
\texttt{quartic\_koszul\_}\allowbreak\texttt{squares} & for the Koszul resolution of $R/(x_{1},\dots,x_{c})$, $c=3,4$, the product $a_{\partial_{k}}a_{\partial_{l}}$ is represented by a nonzero vector exactly when $k\ne l$ and $k,l\le c$, and the coboundaries vanish at the origin, \Cref{lem:freegerm} \\
\texttt{quartic\_line\_}\allowbreak\texttt{two\_planes} & a nonzero vector of $[-6,6]^{4}$ has a nonzero component outside one of two complementary coordinate planes, the step of \Cref{thm:quarticlocal}(iii) \\
\texttt{quartic\_first\_}\allowbreak\texttt{factor\_character} & the pieces of \cite[Example~8.2.4]{Mar25a} have the characters used there, and $\ch(F_{d})=\Theta-\frac d6\Theta^{3}$ with $\chi(F_{d},F_{d})=8d$ for $d\le60$, \Cref{cor:markmanquarticfails}(ii) \\
\texttt{conv\_}\allowbreak\texttt{pure\_}\allowbreak\texttt{weil\_}\allowbreak\texttt{identity} & for $n=2,3,4$ there are $2\cdot4^{n-1}$ vectors $\zeta\in\mu_{4}^{n}$ with $\prod\zeta=\pm1$, half of each sign, and the signed sum of the $e^{c_{1}(L_{\zeta})}$ has coefficient $2\cdot4^{n-1}$ on $\alpha$ and on $\bbar\alpha$ and $0$ on every other monomial in the $c_{j}$, $c'_{j}$, \Cref{prop:weilpieces}(i) \\
\texttt{conv\_}\allowbreak\texttt{cup\_}\allowbreak\texttt{kernel} & at $n=3$ the graphs $\Gamma_{\tau}$ of the six types have $4$ components when $\tau$ has an entry $2$ and $16$ otherwise, so the classes constant on components span $204$ dimensions, $249$ with the $45$ classes of the multiples, out of $525$, \Cref{prop:cupkernel} \\
\texttt{conv\_}\allowbreak\texttt{run\_}\allowbreak\texttt{excess} & for $3\le n\le12$ and every placement of $p$ among the $t_{i}$, a run from $L_{\zeta}$ through distinct multiples back has excess at least $2n-4\ge2$, and a closed walk through the multiples alone at least $2n-3$; at $n=2$ a run of excess $0$ exists, \Cref{prop:nothingenters} \\
\texttt{conv\_}\allowbreak\texttt{esix\_}\allowbreak\texttt{thresholds} & along $M_{2}\to M_{3}\to L_{\zeta}\to M_{1}$ with $\Ext$ degrees $0,6,0$ the shifts are $d,d+1,d-4,d-3$, the product lands in $\Ext$ degree $6$ and $\sigma_{1}=\sigma_{3}=-\sigma_{2}$; $r=24,57,104$ at $n=2,3,4$; $(t_{2}-t_{1})^{6}=1,64,729$, $249-64=185>57$, $249-57=192$, and $(t_{2}-t_{1})^{6}\ge192$ exactly when $t_{2}-t_{1}\ge3$, \Cref{thm:esixdiagonal} \\
\texttt{ff\_}\allowbreak\texttt{partner\_}\allowbreak\texttt{counts} & with $\zeta_{j}=i^{e_{j}}$, every piece has $3$, $6$, $7$ partners of the other parity differing in $1$, $2$, $3$ coordinates; for the $7$ the dimension $D=\prod_{j}|\zeta_{j}-\zeta'_{j}|^{2}$ is $16$ six times and $64$ once, $160$ in all, and the $112$ pairs carry $2560$ classes, \Cref{lem:allthree}, \Cref{prop:explicitconvolution} \\
\texttt{ff\_}\allowbreak\texttt{shift\_}\allowbreak\texttt{table} & in the three placements of the multiples, the shifts $\mu_{i}$, $\nu_{i}$, $\lambda_{i}$, $\kappa_{i}$ are those of the table in the proof, and $\lambda_{i}\le r\le\kappa_{k}$, $\kappa_{i}\ge\lambda_{i}+2$, $\lambda_{i}\le\mu_{i}+2$, $\nu_{i}\le\kappa_{i}+2$, \Cref{lem:allthree} \\
\texttt{ff\_}\allowbreak\texttt{noconvolution\_}\allowbreak\texttt{counts} & the kernels $11$, $19$ of $x_{1}$, $x_{3}$ on a surface at $t=(2,5,6)$, the degrees of the covers, $3^{6}=729$, $16\cdot24^{3}=221184$, $4^{6}=4096$, and the bounds $249-(45+15\cdot9)=69>57$ and $16\cdot12=192>57$, \Cref{thm:noconvolution} \\
\texttt{efour\_}\allowbreak\texttt{partner\_}\allowbreak\texttt{counts} & at $n=4$ the $128$ pieces, $64$ of each parity; every piece has $4$, $12$, $28$, $20$ partners of the other parity differing in $1$, $2$, $3$, $4$ coordinates, those differing in one coordinate differ there by $2$; for the $28$, $D$ is $16$ twenty-four times and $64$ four times, $640$ in all, and $64\cdot640=40960>104$, \Cref{lem:efourshifts}, \Cref{thm:efourtwolevels} \\
\texttt{efour\_}\allowbreak\texttt{run\_}\allowbreak\texttt{drop} & with $p$ at any rank among the four values, every run from $p$ through distinct multiples before and after a marked step, a component or a class of $H^{1}(a,a)$, back to $p$ lowers the shift by at least $4$, \Cref{lem:efourshifts}(iii) \\
\texttt{efour\_}\allowbreak\texttt{two\_}\allowbreak\texttt{level\_}\allowbreak\texttt{degrees} & with $p$ at any rank, for chains $c\to\dots\to W$ and $V\to\dots\to e$ through distinct multiples, not both empty, the degree $4-U+7D$ is not $3$ if $c=e$, not $0$ if $e$ lies above $c$ and not $8$ if below, \Cref{thm:efourtwolevels}(ii) \\
\texttt{vandermonde\_}\allowbreak\texttt{counts} & the genus of the generalised Fermat curve by the Riemann--Hurwitz formula and by adjunction for $2\le d\le7$, $2\le N\le8$; $r!\,d^{N(r-1)}=16,32,54,1536,162$ at the five fibre counts; the units of $\ZZ/d$ for $d=2,3,4,6$; the $70$, $490$, $6125$ balanced orbits in $\PP^{6}$ and the bounds $141$, $981$, $12251$, \Cref{thm:vandermonde}, \Cref{cor:twodiagonal} \\
\texttt{vg\_}\allowbreak\texttt{chain\_}\allowbreak\texttt{determinants} & the intersection matrix of a chain of $m$ vanishing cycles has determinant $1$ for $m$ even and $0$ for $m$ odd, $m\le7$, \Cref{thm:vgvandermonde}(i) \\
\texttt{vg\_}\allowbreak\texttt{quadric\_}\allowbreak\texttt{counts} & $1+\sum_{j>r/2}\binom{N+1}{2j}$ is one plus the number of subsets of even size at least $r+2$ ($N\le9$), $N+2$ for two quadrics and $2$ for a quadric; $2,8,30,94,257$ for $r=4$, $N=5,\dots,9$, \Cref{thm:vgvandermonde}(iii) \\
\texttt{vg\_}\allowbreak\texttt{cubic\_}\allowbreak\texttt{counts} & $1+\binom{N+1}{r+2}\binom{r+2}{r/2+1}$ is one plus the number of balanced characters ($N\le7$, $r=2,4$); $7$, $21$, $71$ for the Fermat cubics, $141$ and $631$ for two cubics in $\PP^{6}$ and $\PP^{8}$, \Cref{thm:vgvandermonde}(iii) \\
\texttt{vg\_}\allowbreak\texttt{triple\_}\allowbreak\texttt{signatures} & for $n_{1},n_{2}<20$ branch points of exponents $1,2$, the signature $(p,q)$ satisfies $n_{1}=2p-q+1$, $n_{2}=2q-p+1$, balanced exactly when $n_{1}=n_{2}$, \Cref{thm:vgvandermonde}(ii) \\
\texttt{cyclic\_}\allowbreak\texttt{base\_}\allowbreak\texttt{count} & the vectors of nonzero residues modulo $m=3,4,6$ with $5$ or $6$ coordinates, sum $0$, order $m$ and $p,q\ge1$ number $38$ up to sign, the cases with five or six branch points of \Cref{prop:cyclicmonodromy}, which the paper proves by hand \\
\texttt{cyclic\_}\allowbreak\texttt{merge\_}\allowbreak\texttt{small} & every such vector with $k$ coordinates, $7\le k\le8$, $11$, $10$ for $m=3,4,6$, has two entries $u,w$ with $u+w\ne0$ whose merge keeps the order $m$ and $p,q\ge1$, a check of \Cref{lem:merge} \\
\texttt{cyclic\_}\allowbreak\texttt{norms} & $x^{2}+xy+y^{2}$ is never $2$ modulo $4$ and is even only for $x,y$ even, so $2$ is not a norm from $\QQ(\sqrt{-3})$; $2$, $3$ and $4$ are norms where they are used, \Cref{lem:newdisc} \\
\texttt{cyclic\_}\allowbreak\texttt{vg\_}\allowbreak\texttt{counts} & $T_{d}(k)$ by a recursion agrees with enumeration; $T_{4}(6)=141$, $T_{6}(6)=1751$, $T_{4}(8)=1107$, $T_{6}(8)=38165$; the formula agrees with a direct count of the characters and gives $142$, $988$, $3950$, $1108$, $1752$, $12258$, $38166$, \Cref{thm:vgvandermonde}(iii) \\
\texttt{cyclic\_}\allowbreak\texttt{nonsplit\_}\allowbreak\texttt{example} & $(1,1,2,2,3,5,5,5)$ modulo $6$ has sum $0$, order $6$, $p=q=3$ and one coordinate $3$, \Cref{rem:vandermonde} \\
\texttt{efour\_}\allowbreak\texttt{three\_}\allowbreak\texttt{targets} & every piece on $E_{0}^{8}$ has $18$, $24$, $21$ partners of its parity differing in $2$, $3$, $4$ coordinates, with groups $\cH^{5}$ of dimensions $64$ (six times) and $16$ (twelve times), $16$, and $0$, adding up to $960$; $32\cdot960=30720$ and $40960-30720=10240$, \Cref{thm:efourthreelevels}(iii) \\
\texttt{efour\_}\allowbreak\texttt{three\_}\allowbreak\texttt{degrees} & with $p$ at any rank and $\delta=1,2$, two runs through distinct multiples, not both empty, give $\Ext$ degree $3+\delta-U+7D$, never $3$ when $c=e$, $0$ when $e$ is above $c$ or $8$ when $e$ is below $c$, \Cref{thm:efourthreelevels}(ii) \\
\texttt{efour\_}\allowbreak\texttt{three\_}\allowbreak\texttt{spectrum} & the character sums of the Cayley graph of the targets are real, with values $960$, $384$, $192$, $96$, $-32$, $-64$, $-128$, $-192$; $64(960+192)/4=18432$, attained by $\zeta_{3}\zeta_{4}\in\{1,i\}$ in both parities, and $40960-18432=22528$, \Cref{thm:efourthreelevels}(iii) \\
\texttt{efour\_}\allowbreak\texttt{spread\_}\allowbreak\texttt{ratios} & the ratios of sum $0$ modulo $4$ that move every coordinate are $21$; the $13$ that move one coordinate by $-1$ have $D=256$ once and $64$ twelve times, $1024$ in all, and the other $8$ have $D=16$, \Cref{lem:efourspread}(ii) \\
\texttt{efour\_}\allowbreak\texttt{spread\_}\allowbreak\texttt{parity} & for each of the $13$ and every set $K$ of coordinates over which the exponents have even sum, that sum is not $2(|K|-1)$ modulo $4$; each of the other $8$ attains it for some $K$, \Cref{lem:efourspread}(ii) \\
\texttt{efour\_}\allowbreak\texttt{cut\_}\allowbreak\texttt{bound} & the $13$ ratios generate the $64$ ratios of sum $0$ in three steps; $|H|(1024-256-64(|H|-2))$ is $1536$, $2560$, $3072$ for $|H|=2,4,8$, and $64|H|\ge1024$ for $|H|=16,32$, \Cref{lem:cayleycut} \\
\texttt{efour\_}\allowbreak\texttt{spread\_}\allowbreak\texttt{three\_}\allowbreak\texttt{counts} & between every piece and the $64$ of the other parity the groups $\cH^{5}$ add up to $1072$ and the groups $\cH^{7}$ to $16$; $64\cdot1072-128\cdot8=67584$ and $64\cdot16=1024$, \Cref{thm:efourspreadthree} \\
\texttt{efour\_}\allowbreak\texttt{five\_}\allowbreak\texttt{values} & with one parity at values in $\{0,2,4\}$ and the other in $\{1,3\}$, every pair of occupied sets is covered by one of the four exclusions; within six values exactly two pairs are left, \Cref{cor:efourfive} \\
\texttt{efour\_}\allowbreak\texttt{cup\_}\allowbreak\texttt{kernel\_}\allowbreak\texttt{counts} & $4\cdot4+6\cdot16\cdot4=400$, $400+3\cdot28=484$, $131\cdot28=3668$, $3584-400=3184$, $484-104=380$; the pieces take $4$ values in each coordinate and $16$ in each pair, \Cref{prop:efourcupkernel} \\
\texttt{efour\_}\allowbreak\texttt{run\_}\allowbreak\texttt{drops} & a run from a piece through one, two or three distinct multiples to a piece lowers the shift by $4$, $5$, $6$, $12$, $13$ or $20$, for every rank of $p$, and never by $7$, \Cref{lem:efourlonely} \\
\texttt{efour\_}\allowbreak\texttt{lonely\_}\allowbreak\texttt{degrees} & for odd $g$ from $5$ to $31$, $X\in\{0,g\}$ and runs with $U$ steps up and $D$ down, $1\le U+D\le3$, the degree $3+X-U+7D$ is $0$ or $8$ only if $X=0$ and $U+D=3$, or $X=7$, $D=0$ and $U=2$; with no run it is at least $8$, \Cref{lem:efourlonely} \\
\texttt{efour\_}\allowbreak\texttt{lonely\_}\allowbreak\texttt{residues} & along three steps of $+1$ or $-7$ the four shifts are pairwise incongruent modulo $4$; for every rank of $p$ the four paths give four patterns of shifts of the multiples, none shared by two shifts at an odd distance below $40$, \Cref{prop:efoursingle} \\
\texttt{efour\_}\allowbreak\texttt{interleaved\_}\allowbreak\texttt{cut} & the $28$ ratios that move three coordinates and change the parity weigh $640$, $16$ twenty-four times and $64$ four times, and generate the $128$ ratios in three steps; the bounds $1152$, $1792$, $2688$, $3328$, $1024$, $1024$ for $|H|=2,4,\dots,64$ exceed $640$, \Cref{prop:efourpairs} \\
\texttt{efour\_}\allowbreak\texttt{six\_}\allowbreak\texttt{values} & with one parity at values in $\{0,2,4\}$ and the other in $\{1,3,5\}$, every pair of occupied sets is covered by the exclusions of \Cref{cor:efourfive}, by single shifts at an odd distance or by the interleaved sets; single shifts at every odd distance up to $41$ are covered; $64\cdot28=1792>104$, \Cref{cor:efoursix} \\
\texttt{divisor\_}\allowbreak\texttt{profiles} & the Betti numbers of the connected sum of two real $n$-tori have Euler characteristic $0$ for odd $n$ and $-2$ for even $n$, $2\le n\le40$; at $n=5$ they are $1,10,20,20,10,1$; for $5\le n\le16$, eight values of $d$ and $0\le a\le3$, $|b|\le3$ the profile of part (i) meets every condition, \Cref{prop:divisortemplate}(i) \\
\texttt{divisor\_}\allowbreak\texttt{endo\_}\allowbreak\texttt{positive} & $r^{2}(20d+1)+6<32r^{2}d^{2}$ for all $d,r\ge1$, so $6\chi(D,\End_{0}G)>0$, \Cref{prop:divisortemplate}(iv) \\
\texttt{divisor\_}\allowbreak\texttt{rank\_}\allowbreak\texttt{two\_}\allowbreak\texttt{mod3} & $3$ divides $2(d+1)$ exactly when $d\equiv2\pmod3$, for every $d$, \Cref{prop:divisortemplate}(ii) \\
\texttt{hh\_}\allowbreak\texttt{annihilator\_}\allowbreak\texttt{small} & for $n=1,2,3$ and every monomial $S$ of $HH^{*}$: if $S$ meets $P$ and $Q$ it kills $\alpha_{+}$ and $\alpha_{-}$; if it lies in $Q$ (in $P$) it sends $\alpha_{+}$ ($\alpha_{-}$) to a monomial and kills the other; the images are distinct except for $S=Q$, $S'=P$ in degree $2n$; so $\dim\operatorname{Ann}_{HH^{k}}(\omega)=\binom{4n}{k}-2\binom{2n}{k}+\delta_{k,2n}$, the mixed products of two wedges number $n^{2}$, and $2^{2n+1}-1$ monomials lie outside the ideal of $P\wedge Q$, \Cref{lem:annihilator}, \Cref{thm:hhsplitting}, \Cref{cor:hhfactor} \\
\texttt{polarised\_}\allowbreak\texttt{number\_}\allowbreak\texttt{three} & at $n=3$, for two characters of each Hankel rank $\rho=0,1,2,3$, each with $u=1$ and $u=2+3i$, $\dim\operatorname{Ann}_{HT^{2}}(\gamma)=n^{2}(4-\rho)$ and the contraction has rank $(4+\rho)n^{2}-2n$: exact annihilating vectors, and images independent modulo $998244353$, \Cref{thm:p2primenumber} \\
\texttt{polarised\_}\allowbreak\texttt{number\_}\allowbreak\texttt{four\_}\allowbreak\texttt{rho0} & the same at $n=4$ for $\gamma=(2+3i)\alpha_{+}+(2-3i)\alpha_{-}$: annihilator $64$, rank $56$, \Cref{thm:p2primenumber} \\
\texttt{polarised\_}\allowbreak\texttt{number\_}\allowbreak\texttt{four\_}\allowbreak\texttt{rho1} & the same for $e^{\theta}$ with that Weil part: annihilator $48$, rank $72$, \Cref{thm:p2primenumber} \\
\texttt{polarised\_}\allowbreak\texttt{number\_}\allowbreak\texttt{four\_}\allowbreak\texttt{rho2} & the same for $\theta$ with that Weil part: annihilator $32$, rank $88$, \Cref{thm:p2primenumber} \\
\texttt{polarised\_}\allowbreak\texttt{number\_}\allowbreak\texttt{four\_}\allowbreak\texttt{rho3} & the same for a general integral polynomial part: annihilator $16$, rank $104$, \Cref{thm:p2primenumber} \\
\texttt{weil\_}\allowbreak\texttt{product\_}\allowbreak\texttt{formula} & the coefficient $c_{n}(t)$ of $w^{(T)}$, $|T|=t$, in $\omega_{1}+\sqrt{-d}\,\omega_{2}$ satisfies $d^{n}c_{n}(t)=d^{2n-t}(\sqrt{-d})^{t}$, seven fields, $n\le10$, \eqref{eq:omegaprod} \\
\texttt{weil\_}\allowbreak\texttt{multiplicative\_}\allowbreak\texttt{coefficients} & $c_{n_{1}}(t_{1})\,c_{n_{2}}(t_{2})=c_{n_{1}+n_{2}}(t_{1}+t_{2})$ in $\ZZ[\sqrt{-d}]$, seven fields, $n_{1},n_{2}\le6$, \eqref{eq:weilmult} \\
\texttt{weil\_}\allowbreak\texttt{multiplicative\_}\allowbreak\texttt{exterior} & in the exterior algebra over $\ZZ[\sqrt{-d}]$, $\prod_{j}(d\,x_{j}+\sqrt{-d}\,y_{j})=d^{n}\omega$ for $n\le3$, and for every composition $n=n_{1}+\dots+n_{k}$ the product of the classes $\omega$ of the factors is the class of the product, computed with signs, for $n\le3$ and five fields and for $n=4$ and $d=1,3$, \eqref{eq:omegaprod}, \Cref{thm:weilmult} \\
\texttt{split\_}\allowbreak\texttt{member\_}\allowbreak\texttt{closed\_}\allowbreak\texttt{form} & on $X\times\widehat{X}$, $n!\,\omega=(-1)^{n(n-1)/2}(\gamma+\sqrt{-d}\,\ell)^{n}$ for $n\le3$; $2\omega_{1}=d\ell^{2}-\gamma^{2}$ and $\omega_{2}=-\gamma\ell$ at $n=2$; $6\omega_{1}=3d\gamma\ell^{2}-\gamma^{3}$ and $6\omega_{2}=d\ell^{3}-3\gamma^{2}\ell$ at $n=3$; five fields, \Cref{ex:splitsmall} \\
\texttt{secant\_}\allowbreak\texttt{chern\_}\allowbreak\texttt{recursion} & the recursion $\gamma_{k+1}=b\gamma_{k}+dk(a-k+1)\gamma_{k-1}$ agrees with Newton's identities for $k\le8$, eight squarefree $d\le11$, $0\le a\le4$ and $|b|\le4$, a box on which agreement is the identity, \Cref{prop:notforbidden}(ii) \\
\texttt{secant\_}\allowbreak\texttt{discriminant} & $a\gamma_{2}-(a-1)b^{2}=a^{2}d+b^{2}$, which is positive unless $a=b=0$, \Cref{prop:notforbidden}(iii) \\
\texttt{secant\_}\allowbreak\texttt{rank\_}\allowbreak\texttt{one\_}\allowbreak\texttt{chern} & at $a=1$, $\gamma_{2}=b^{2}+d$, $\gamma_{3}=b(b^{2}+d)$ and $\gamma_{4}=(b^{2}+d)(b^{2}-3d)$, the Chern classes of $\cI_{Z}(b\Theta)$ in \Cref{prop:lowrankburch} \\
\texttt{weil\_}\allowbreak\texttt{plane\_}\allowbreak\texttt{gram} & $\eta\,\omega_{1}=\eta\,\omega_{2}=0$, $\eta^{2n}$ is $(2n)!$ times the top monomial, and the Gram matrix of the rational Weil plane is $(-1)^{n}2^{2n-1}\operatorname{diag}(d^{n},d^{n-1})$, for $n\le3$ and five fields, \Cref{thm:p2endo} \\
\texttt{lattice\_}\allowbreak\texttt{secant\_}\allowbreak\texttt{moments} & the binomial moments of $a\,u_{\Theta}+b\,v_{\Theta}$ are those of the theorem, on a box on which agreement is the identity, \Cref{lem:binomialmoments}, \Cref{thm:lattice} \\
\texttt{lattice\_}\allowbreak\texttt{u3v\_}\allowbreak\texttt{moments} & $24\,m(u_{\Theta}+3v_{\Theta})=\bigl(24,72,-12(d+3),24,(d+9)(d-2)\bigr)$, \Cref{thm:lattice} \\
\texttt{lattice\_}\allowbreak\texttt{u3v\_}\allowbreak\texttt{criterion} & for every $d$, $d$ is odd and $24\mid(d+9)(d-2)$ exactly when $d\equiv15,23\pmod{24}$, \Cref{thm:lattice} \\
\texttt{lattice\_}\allowbreak\texttt{moments\_}\allowbreak\texttt{d3} & at $d=3$, $m(u_{\Theta}+v_{\Theta})=(1,1,-2,\tfrac43,-\tfrac12)$, with least common denominator $6$, \Cref{thm:lattice} \\
\texttt{lattice\_}\allowbreak\texttt{smooth\_}\allowbreak\texttt{support} & $m(\cO_{S})=\bigl(0,0,N,-3N,N(83-2N)/12\bigr)$, not integral for $N=5,\dots,8$, \Cref{rem:latticedefect} \\
\texttt{burch\_}\allowbreak\texttt{rank\_}\allowbreak\texttt{one\_}\allowbreak\texttt{c4} & at $r=1$, $c_{1}(E)=-\frac b3\Theta$ kills $c_{3}(G)$, and then $72\,c_{4}(G)=-(b^{2}+d)(b^{2}+9d)\ne0$, \Cref{prop:lowrankburch}(i) \\
\texttt{burch\_}\allowbreak\texttt{rank\_}\allowbreak\texttt{two\_}\allowbreak\texttt{c4} & at $r=2$, $24\,c_{4}(G)=\gamma_{2}(b^{2}-3d+4ab+6m)$, and $\chi(E)=a^{4}-2a^{2}m+\frac12m^{2}$ by Newton's identities, an integer exactly when $m$ is even, \Cref{prop:lowrankburch}(ii) \\
\texttt{burch\_}\allowbreak\texttt{rank\_}\allowbreak\texttt{two\_}\allowbreak\texttt{parity} & for every $d$, $12\mid3d+3$ exactly when $d\equiv3\pmod4$; as $3d-9-12a=(3d+3)-12(a+1)$ this is the case $b=3$, \Cref{prop:lowrankburch}(ii) \\
\texttt{jacobian\_}\allowbreak\texttt{ring\_}\allowbreak\texttt{counts} & the coefficients of $(1+s+\dots+s^{d-2})^{n+2}$ agree with inclusion and exclusion, are symmetric and add up to $(d-1)^{n+2}$, and the Euler number of a hypersurface agrees with its two closed forms, for $n\le8$, $2\le d\le8$, \Cref{prop:f3primefibration} \\
\texttt{sextic\_}\allowbreak\texttt{fourfold\_}\allowbreak\texttt{hodge} & the sextic fourfold has primitive Hodge numbers $1,426,1751,426,1$, Euler number $2610$ and $b_{4}=2606$; $h^{4,0}=\binom{d-1}{5}$ for $3\le d\le7$, zero exactly in the Fano degrees, \Cref{rem:f3primesharp} \\
\texttt{blowup\_}\allowbreak\texttt{sextic\_}\allowbreak\texttt{hodge} & $\mathrm{Bl}_{Y}\PP^{6}$, for a sextic fourfold $Y$ in a hyperplane, has $(h^{6,0},h^{5,1},h^{4,2},h^{3,3})=(0,1,426,1753)$, Hodge symmetry, only even cohomology, $h^{p,0}=0$ for $p>0$ and total Betti number $2617=e(\PP^{6})+e(Y)$, \Cref{rem:f3primesharp} \\
\texttt{blowup\_}\allowbreak\texttt{retrieval} & on a $\PP^{1}$-bundle over a base with $H^{*}(Y)=\QQ[y]/(y^{5})$, $-\rho_{*}j^{*}j_{*}\rho^{*}a=a$ in every degree and $\xi^{2}=-c_{1}\xi-c_{2}$, for $c_{1}=sy$, $c_{2}=ty^{2}$, $|s|,|t|\le6$, \Cref{prop:f3primefibration} \\
\texttt{zero\_}\allowbreak\texttt{cycle\_}\allowbreak\texttt{degrees} & a class of degree $2p$ on an $n$-fold needs an input only for $2\le p\le n/2$; on a blow-up of a rational sixfold only in degree four on a fourfold centre, and on a uniruled fivefold only in degree four on a fourfold, \Cref{prop:f3primefibration}, \Cref{rem:f3primesharp} \\
\texttt{rigidity\_}\allowbreak\texttt{polarisation} & for $n=2,3,4$ the contraction map on $H^{1,1}_{\RR}$, of dimension $4n^{2}$, has rank $4n^{2}-1$ and kills the polarisation, \Cref{prop:rigidity} \\
\texttt{rigidity\_}\allowbreak\texttt{norm\_}\allowbreak\texttt{kernels} & a norm-character class of rank $r$ has kernel of dimension $n(n-r)$ on the tangent space, for every $r\le n$, $n=2,3,4$, two shapes each, and a class without structure has kernel zero, \Cref{prop:rigidity} \\
\texttt{object\_}\allowbreak\texttt{size\_}\allowbreak\texttt{bound} & $\alpha_{\pm}^{2}=0$, $\int\omega_{1}^{2}=2$, and $\chi(E,E)=(-1)^{n}2N^{2}$ for $\ch(E)=r+N\omega_{1}$ whatever $r$, for $n\le4$, \Cref{prop:objectsize} \\
\texttt{tangent\_}\allowbreak\texttt{is\_}\allowbreak\texttt{annihilator\_}\allowbreak\texttt{two}, \texttt{tangent\_}\allowbreak\texttt{is\_}\allowbreak\texttt{annihilator\_}\allowbreak\texttt{three}, \texttt{tangent\_}\allowbreak\texttt{is\_}\allowbreak\texttt{annihilator\_}\allowbreak\texttt{three'}, \texttt{tangent\_}\allowbreak\texttt{is\_}\allowbreak\texttt{annihilator\_}\allowbreak\texttt{four} & in a second model, for $n=2,3,4$ and several planes: the Weil classes are real of type $(n,n)$, their annihilator in $H^{0,2}$ has dimension $n^{2}$, and $\Hom(P,Q)\to H^{0,2}$ is injective onto it, \Cref{lem:annihilator}, \Cref{thm:anntangent} \\
\texttt{tangent\_}\allowbreak\texttt{deformations\_}\allowbreak\texttt{four}, \texttt{tangent\_}\allowbreak\texttt{hodge\_}\allowbreak\texttt{locus\_}\allowbreak\texttt{four} & at $n=4$ the polarised first-order deformations have dimension $36$, the $K$-linear ones $16$, and those killing the Weil line are exactly the $K$-linear ones, \Cref{prop:firstorder} \\
\texttt{weil\_}\allowbreak\texttt{tori\_}\allowbreak\texttt{types\_}\allowbreak\texttt{two}, \texttt{weil\_}\allowbreak\texttt{tori\_}\allowbreak\texttt{types\_}\allowbreak\texttt{three}, \texttt{weil\_}\allowbreak\texttt{tori\_}\allowbreak\texttt{kappa\_}\allowbreak\texttt{three} & the Weil classes are of type $(n,n)$ on the whole $K$-linear family, the polarisation only at the balanced member, and $\alpha_{+}\kappa^{n}=0$ for $\kappa$ in $V_{+}\otimes V_{-}$, at $n=2$ ($d=1,3$) and $n=3$ ($d=2$), \Cref{setup:weiltori} \\
\texttt{weil\_}\allowbreak\texttt{tori\_}\allowbreak\texttt{hodge\_}\allowbreak\texttt{two}, \texttt{weil\_}\allowbreak\texttt{tori\_}\allowbreak\texttt{orbit} & a very general Weil torus of dimension four (and six, through the torus orbit of one member) has no Hodge classes in degree $2k$, $0<k<n$, and only the Weil plane in degree $2n$, \Cref{lem:weiltorushodge}, \Cref{lem:weiltorussubsets} \\
\texttt{hochschild\_}\allowbreak\texttt{profile\_}\allowbreak\texttt{two} & the Hochschild profile at $n=2$ on the eight shapes, each with $u=1$ and $u=2+3i$: $(1,8,\rho_{2},8,1)$ with $\rho_{2}=12,16,20,24$ as the Hankel rank is $0,1,2,3$, \Cref{prop:p2primeprofile} \\
\texttt{hochschild\_}\allowbreak\texttt{profile\_}\allowbreak\texttt{three\_}\allowbreak\texttt{pure}, \texttt{hochschild\_}\allowbreak\texttt{profile\_}\allowbreak\texttt{three\_}\allowbreak\texttt{cn}, \texttt{hochschild\_}\allowbreak\texttt{profile\_}\allowbreak\texttt{three\_}\allowbreak\texttt{generic\_}\allowbreak\texttt{low}, \texttt{hochschild\_}\allowbreak\texttt{profile\_}\allowbreak\texttt{three\_}\allowbreak\texttt{generic\_}\allowbreak\texttt{middle}, \texttt{hochschild\_}\allowbreak\texttt{profile\_}\allowbreak\texttt{three\_}\allowbreak\texttt{generic\_}\allowbreak\texttt{high} & the profiles at $n=3$ of the pure character, of $c_{3}\theta^{3}$ and of a general shape, $(1,12,57,112,57,12,1)$ for the last, each rank exact by independent images and exact relations, \Cref{prop:p2primeprofile} \\
\texttt{hochschild\_}\allowbreak\texttt{profile\_}\allowbreak\texttt{special\_}\allowbreak\texttt{two}, \texttt{hochschild\_}\allowbreak\texttt{profile\_}\allowbreak\texttt{special\_}\allowbreak\texttt{three}, \texttt{hochschild\_}\allowbreak\texttt{profile\_}\allowbreak\texttt{special\_}\allowbreak\texttt{three\_}\allowbreak\texttt{second} & for $c_{n}\theta^{n}$ alone the middle rank drops exactly at $|u|=\binom{n}{a}n!\,|c_{n}|$ and only in degree $n$: $\rho_{2}=22,23$ at $n=2$, $\rho_{3}=110$ at $n=3$, \Cref{prop:p2primeprofile} \\
\texttt{descent\_}\allowbreak\texttt{weil\_}\allowbreak\texttt{quadratic}, \texttt{descent\_}\allowbreak\texttt{weil\_}\allowbreak\texttt{quartic}, \texttt{descent\_}\allowbreak\texttt{weil\_}\allowbreak\texttt{sextic\_}\allowbreak\texttt{one}, \texttt{descent\_}\allowbreak\texttt{weil\_}\allowbreak\texttt{sextic\_}\allowbreak\texttt{two} & the map $P$ carries $W_{F}(B\times Y)$ isomorphically onto $W_{F}(B)$ and kills it when the Weil class of $Y$ is replaced by its polarisation, for four imaginary quadratic fields, $\QQ(\zeta_{5})$, $\QQ(\zeta_{8})$ and $\QQ(\zeta_{7})$, \Cref{prop:descent} \\
\texttt{descent\_}\allowbreak\texttt{discriminant} & $H_{B}+\operatorname{diag}(1,-t)$ has signature $(n+1,n+1)$ and $t$ times the discriminant of $H_{B}$, for $n\le5$ and five values of $t$, \Cref{prop:descent} \\
\texttt{descent\_}\allowbreak\texttt{scalar\_}\allowbreak\texttt{extension} & $j^{*}W_{F}(B_{F})=W_{K}(B)$ for $\QQ(i)\subset\QQ(\zeta_{8})$, $\QQ(\sqrt{-7})\subset\QQ(\zeta_{7})$ and $\QQ(\sqrt{-3})\subset\QQ(\zeta_{9})$, $n=1,2$, \Cref{prop:descent} \\
\texttt{wirtinger\_}\allowbreak\texttt{volumes}, \texttt{wirtinger\_}\allowbreak\texttt{primitive}, \texttt{wirtinger\_}\allowbreak\texttt{pairings} & on a Mumford square, $\theta^{4}=24\,\mathrm{vol}$, $\theta_{Y}^{8}=8!\,\mathrm{vol}$, the decomposition $\wedge^{2}V=\CC\theta+U_{12}+U_{13}+U_{23}$ into primitive pieces, and the pairings $\int\pi_{0}\theta_{Y}^{6}=2880$, $\int\pi_{ij}\theta_{Y}^{6}=0$, $L(\theta_{Y}^{2})=56$, \Cref{ex:mumfordmass} \\
\texttt{mumford\_}\allowbreak\texttt{motivic\_}\allowbreak\texttt{sp}, \texttt{mumford\_}\allowbreak\texttt{motivic\_}\allowbreak\texttt{irreducible} & $\mathfrak{g}=\mathfrak{sl}_{2}^{3}$ and the $27$ products span $\mathfrak{sp}(V,\psi)$, of dimension $36$; $P$ is irreducible, and $S^{2}V$ is the sum of four irreducible summands, \Cref{prop:mumfordmotivic} \\
\texttt{mumford\_}\allowbreak\texttt{motivic\_}\allowbreak\texttt{permutations}, \texttt{mumford\_}\allowbreak\texttt{motivic\_}\allowbreak\texttt{commutant} & the six permutations of the factors lie in $\Sp(V,\psi)$ and permute the factors of $\mathfrak{g}$; the commutant of $\mathfrak{g}$ is the scalars; the subgroups of $S_{3}$ stable under $A_{3}$ are $1$, $A_{3}$, $S_{3}$, \Cref{prop:mumfordmotivic} \\
\texttt{mumford\_}\allowbreak\texttt{motivic\_}\allowbreak\texttt{projectors}, \texttt{mumford\_}\allowbreak\texttt{motivic\_}\allowbreak\texttt{cycle} & in the model of item (LX) the projectors $\pi_{0},\pi_{12},\pi_{13},\pi_{23}$ are idempotents of traces $1,9,9,9$, hence of these ranks, and $\zeta$ fixes $\pi_{0}$ and carries $\pi_{12}$ to $\pi_{23}$, $\pi_{23}$ to $\pi_{13}$, $\pi_{13}$ to $\pi_{12}$, \Cref{setup:omegaa} \\
\texttt{mumford\_}\allowbreak\texttt{motivic\_}\allowbreak\texttt{hyperdeterminant}, \texttt{mumford\_}\allowbreak\texttt{motivic\_}\allowbreak\texttt{weyl} & Cayley's hyperdeterminant is invariant under $G$ and the permutations but not under $\Sp(V,\psi)$; $\dim(S^{4}V)^{G}=1$ and $\dim(S^{4}V)^{\Sp}=0$, \Cref{prop:mumfordmotivic} \\
\texttt{mumford\_}\allowbreak\texttt{motivic\_}\allowbreak\texttt{invariants} & the non-crossing pairings give bases of the invariants, $1,8,125$ products over the factors, with $4,45$ orbits of $A_{3}$ and $4,35$ of $S_{3}$, and the pairings of $\psi$ have Gram matrices of ranks $1,3,15$, \Cref{prop:mumfordformal} \\
\texttt{mumford\_}\allowbreak\texttt{three\_}\allowbreak\texttt{adic\_}\allowbreak\texttt{anisotropic} & $\langle6,-2,-2\rangle$ has no primitive zero modulo $27$, so it is anisotropic over $\QQ_{3}$, \Cref{rem:mumfordks} \\
\texttt{k3\_}\allowbreak\texttt{hilbert\_}\allowbreak\texttt{square\_}\allowbreak\texttt{betti}, \texttt{k3\_}\allowbreak\texttt{fujiki\_}\allowbreak\texttt{determinant}, \texttt{k3\_}\allowbreak\texttt{riemann\_}\allowbreak\texttt{roch}, \texttt{k3\_}\allowbreak\texttt{partition\_}\allowbreak\texttt{sizes} & the Betti numbers $(1,0,23,0,276,0,23,0,1)$ of $S^{[2]}$; the determinant of the Fujiki pairing on $\mathrm{Sym}^{2}H^{2}$, $2^{46}\cdot25$; $\int c_{2}^{2}=828$ and $\chi(L)=\binom{q(L)/2+3}{2}$; and the least $n$ with $t$ distinct part sizes, $t(t+1)/2$, \Cref{prop:k3powers}, \Cref{prop:k3ntype} \\
\texttt{sextic\_}\allowbreak\texttt{lattice\_}\allowbreak\texttt{fields}, \texttt{sextic\_}\allowbreak\texttt{lattice\_}\allowbreak\texttt{euler\_}\allowbreak\texttt{even}, \texttt{sextic\_}\allowbreak\texttt{lattice\_}\allowbreak\texttt{shapes}, \texttt{sextic\_}\allowbreak\texttt{lattice\_}\allowbreak\texttt{exp\_}\allowbreak\texttt{classes} & the sixteen cubic fields of the cases; $\chi(v,v)$ is even; $r^{3}-2r^{2}\le4$ on every support and rank pattern, with equality only at $(54,112)$; and the two classes $\operatorname{Re}e^{i\theta}$, $\operatorname{Im}e^{i\theta}$, \Cref{prop:sexticintegral} \\
\texttt{sextic\_}\allowbreak\texttt{lattice\_}\allowbreak\texttt{case\_}\allowbreak\texttt{00}, \texttt{sextic\_}\allowbreak\texttt{lattice\_}\allowbreak\texttt{case\_}\allowbreak\texttt{01}, \texttt{sextic\_}\allowbreak\texttt{lattice\_}\allowbreak\texttt{case\_}\allowbreak\texttt{02}, \texttt{sextic\_}\allowbreak\texttt{lattice\_}\allowbreak\texttt{case\_}\allowbreak\texttt{03}, \texttt{sextic\_}\allowbreak\texttt{lattice\_}\allowbreak\texttt{case\_}\allowbreak\texttt{04}, \texttt{sextic\_}\allowbreak\texttt{lattice\_}\allowbreak\texttt{case\_}\allowbreak\texttt{05}, \texttt{sextic\_}\allowbreak\texttt{lattice\_}\allowbreak\texttt{case\_}\allowbreak\texttt{06}, \texttt{sextic\_}\allowbreak\texttt{lattice\_}\allowbreak\texttt{case\_}\allowbreak\texttt{07}, \texttt{sextic\_}\allowbreak\texttt{lattice\_}\allowbreak\texttt{case\_}\allowbreak\texttt{08}, \texttt{sextic\_}\allowbreak\texttt{lattice\_}\allowbreak\texttt{case\_}\allowbreak\texttt{09}, \texttt{sextic\_}\allowbreak\texttt{lattice\_}\allowbreak\texttt{case\_}\allowbreak\texttt{10}, \texttt{sextic\_}\allowbreak\texttt{lattice\_}\allowbreak\texttt{case\_}\allowbreak\texttt{11}, \texttt{sextic\_}\allowbreak\texttt{lattice\_}\allowbreak\texttt{case\_}\allowbreak\texttt{12}, \texttt{sextic\_}\allowbreak\texttt{lattice\_}\allowbreak\texttt{case\_}\allowbreak\texttt{13}, \texttt{sextic\_}\allowbreak\texttt{lattice\_}\allowbreak\texttt{case\_}\allowbreak\texttt{14}, \texttt{sextic\_}\allowbreak\texttt{lattice\_}\allowbreak\texttt{case\_}\allowbreak\texttt{15}, \texttt{sextic\_}\allowbreak\texttt{lattice\_}\allowbreak\texttt{case\_}\allowbreak\texttt{16}, \texttt{sextic\_}\allowbreak\texttt{lattice\_}\allowbreak\texttt{case\_}\allowbreak\texttt{17}, \texttt{sextic\_}\allowbreak\texttt{lattice\_}\allowbreak\texttt{case\_}\allowbreak\texttt{18}, \texttt{sextic\_}\allowbreak\texttt{lattice\_}\allowbreak\texttt{case\_}\allowbreak\texttt{19}, \texttt{sextic\_}\allowbreak\texttt{lattice\_}\allowbreak\texttt{case\_}\allowbreak\texttt{20}, \texttt{sextic\_}\allowbreak\texttt{lattice\_}\allowbreak\texttt{case\_}\allowbreak\texttt{21}, \texttt{sextic\_}\allowbreak\texttt{lattice\_}\allowbreak\texttt{case\_}\allowbreak\texttt{22}, \texttt{sextic\_}\allowbreak\texttt{lattice\_}\allowbreak\texttt{case\_}\allowbreak\texttt{23}, \texttt{sextic\_}\allowbreak\texttt{lattice\_}\allowbreak\texttt{case\_}\allowbreak\texttt{24}, \texttt{sextic\_}\allowbreak\texttt{lattice\_}\allowbreak\texttt{case\_}\allowbreak\texttt{25}, \texttt{sextic\_}\allowbreak\texttt{lattice\_}\allowbreak\texttt{case\_}\allowbreak\texttt{26}, \texttt{sextic\_}\allowbreak\texttt{lattice\_}\allowbreak\texttt{case\_}\allowbreak\texttt{27}, \texttt{sextic\_}\allowbreak\texttt{lattice\_}\allowbreak\texttt{case\_}\allowbreak\texttt{28}, \texttt{sextic\_}\allowbreak\texttt{lattice\_}\allowbreak\texttt{case\_}\allowbreak\texttt{29}, \texttt{sextic\_}\allowbreak\texttt{lattice\_}\allowbreak\texttt{case\_}\allowbreak\texttt{30}, \texttt{sextic\_}\allowbreak\texttt{lattice\_}\allowbreak\texttt{case\_}\allowbreak\texttt{31}, \texttt{sextic\_}\allowbreak\texttt{lattice\_}\allowbreak\texttt{case\_}\allowbreak\texttt{32}, \texttt{sextic\_}\allowbreak\texttt{lattice\_}\allowbreak\texttt{case\_}\allowbreak\texttt{33}, \texttt{sextic\_}\allowbreak\texttt{lattice\_}\allowbreak\texttt{case\_}\allowbreak\texttt{34}, \texttt{sextic\_}\allowbreak\texttt{lattice\_}\allowbreak\texttt{case\_}\allowbreak\texttt{35}, \texttt{sextic\_}\allowbreak\texttt{lattice\_}\allowbreak\texttt{case\_}\allowbreak\texttt{36}, \texttt{sextic\_}\allowbreak\texttt{lattice\_}\allowbreak\texttt{case\_}\allowbreak\texttt{37}, \texttt{sextic\_}\allowbreak\texttt{lattice\_}\allowbreak\texttt{case\_}\allowbreak\texttt{38}, \texttt{sextic\_}\allowbreak\texttt{lattice\_}\allowbreak\texttt{case\_}\allowbreak\texttt{39}, \texttt{sextic\_}\allowbreak\texttt{lattice\_}\allowbreak\texttt{case\_}\allowbreak\texttt{40}, \texttt{sextic\_}\allowbreak\texttt{lattice\_}\allowbreak\texttt{case\_}\allowbreak\texttt{41}, \texttt{sextic\_}\allowbreak\texttt{lattice\_}\allowbreak\texttt{case\_}\allowbreak\texttt{42}, \texttt{sextic\_}\allowbreak\texttt{lattice\_}\allowbreak\texttt{case\_}\allowbreak\texttt{43}, \texttt{sextic\_}\allowbreak\texttt{lattice\_}\allowbreak\texttt{case\_}\allowbreak\texttt{44}, \texttt{sextic\_}\allowbreak\texttt{lattice\_}\allowbreak\texttt{case\_}\allowbreak\texttt{45}, \texttt{sextic\_}\allowbreak\texttt{lattice\_}\allowbreak\texttt{case\_}\allowbreak\texttt{46}, \texttt{sextic\_}\allowbreak\texttt{lattice\_}\allowbreak\texttt{case\_}\allowbreak\texttt{47} & for each of the $48$ pairs of a field and $q$: the integral classes of $S(0,q)$ are the span of the recorded basis, with the recorded Gram matrix of $-\chi/8$, \Cref{prop:sexticintegral} \\
\texttt{sextic\_}\allowbreak\texttt{weil\_}\allowbreak\texttt{units}, \texttt{sextic\_}\allowbreak\texttt{weil\_}\allowbreak\texttt{imaginary}, \texttt{sextic\_}\allowbreak\texttt{weil\_}\allowbreak\texttt{shape\_}\allowbreak\texttt{classes}, \texttt{sextic\_}\allowbreak\texttt{weil\_}\allowbreak\texttt{least\_}\allowbreak\texttt{values} & the units with $2+a=\beta^{2}$; the twenty cases in which $F$ contains an imaginary quadratic field; the classes of least norm with an $F$-Weil part; and the least values of $-\chi$, $192$ only over $\QQ(\zeta_{7})^{+}$ at $q=3+a$, \Cref{lem:sexticweil}; the least values are not used in the paper \\
\texttt{sextic\_}\allowbreak\texttt{weil\_}\allowbreak\texttt{scan\_}\allowbreak\texttt{0}, \texttt{sextic\_}\allowbreak\texttt{weil\_}\allowbreak\texttt{scan\_}\allowbreak\texttt{1}, \texttt{sextic\_}\allowbreak\texttt{weil\_}\allowbreak\texttt{scan\_}\allowbreak\texttt{2}, \texttt{sextic\_}\allowbreak\texttt{weil\_}\allowbreak\texttt{scan\_}\allowbreak\texttt{3}, \texttt{sextic\_}\allowbreak\texttt{weil\_}\allowbreak\texttt{scan\_}\allowbreak\texttt{4}, \texttt{sextic\_}\allowbreak\texttt{weil\_}\allowbreak\texttt{scan\_}\allowbreak\texttt{5}, \texttt{sextic\_}\allowbreak\texttt{weil\_}\allowbreak\texttt{scan\_}\allowbreak\texttt{6}, \texttt{sextic\_}\allowbreak\texttt{weil\_}\allowbreak\texttt{scan\_}\allowbreak\texttt{7} & the scans of the lattices by fraction-free Schur complements: no integral class with $-\chi\le24$, and the least $-\chi$ with an $F$-Weil part as recorded, beyond \Cref{prop:sexticintegral} and not used in the paper \\
\texttt{fm\_}\allowbreak\texttt{powers\_}\allowbreak\texttt{computed} & $\Phi_{\cP}(\theta^{\prime k})=(-1)^{g(g+1)/2+k}\frac{k!}{(g-k)!}\theta^{g-k}$, computed in the exterior algebra of $V\oplus V^{\vee}$ for $g\le4$, \Cref{thm:fmtheta} \\
\texttt{mumford\_}\allowbreak\texttt{routes\_}\allowbreak\texttt{target\_}\allowbreak\texttt{open}, \texttt{mumford\_}\allowbreak\texttt{routes\_}\allowbreak\texttt{equivalent}, \texttt{mumford\_}\allowbreak\texttt{routes\_}\allowbreak\texttt{minimal}, \texttt{mumford\_}\allowbreak\texttt{routes\_}\allowbreak\texttt{stronger}, \texttt{mumford\_}\allowbreak\texttt{routes\_}\allowbreak\texttt{conjecture} & the Horn system of the routes to the Mumford target: it is closed and the target is not proved; its equivalent forms; the twelve minimal routes; the stronger routes; and the minimal sets $\{\mathrm{HC}\}$, $\{F_{3}\}$, $\{L,M\}$ yielding the Hodge conjecture, \Cref{cor:mumfordone}, \Cref{prop:f3ishc} \\
\texttt{targets\_}\allowbreak\texttt{invariant\_}\allowbreak\texttt{dims}, \texttt{targets\_}\allowbreak\texttt{invariant\_}\allowbreak\texttt{ring} & the invariants of $\mathfrak{sl}_{2}^{3}$ on $X\times X$ have dimensions $1,0,3,0,8,0,16,0,28,\dots$, and are generated in degrees two and four with two exceptional classes, \Cref{thm:mumfordcm} \\
\texttt{targets\_}\allowbreak\texttt{cm\_}\allowbreak\texttt{counts}, \texttt{targets\_}\allowbreak\texttt{cm\_}\allowbreak\texttt{generation}, \texttt{targets\_}\allowbreak\texttt{two\_}\allowbreak\texttt{branches} & at a CM point the Hodge classes number $1,0,16,0,132,\dots$ on $X_{c}\times X_{c}$ and are generated by divisor classes and pull-backs; the two branches of classes killed by $Q\,HH^{1}$ and $P\,HH^{1}$, \Cref{thm:mumfordcm}, \Cref{cor:mumfordequiv} \\
\texttt{lefschetz\_}\allowbreak\texttt{pontryagin}, \texttt{lefschetz\_}\allowbreak\texttt{family\_}\allowbreak\texttt{operator} & the Lefschetz operator is a Pontryagin product on ten types of polarisation, and $[L,\Lambda]=H$ for the operator of the theorem on four product families, \Cref{prop:pontryaginlambda}, \Cref{thm:lefschetzfamily} \\
\texttt{split\_}\allowbreak\texttt{resolution\_}\allowbreak\texttt{data} & the scalar $(e^{2}+d)/2$ is positive, the vanishing degrees of $H^{*}(\cO(m\Theta))$ exclude $2$, the three resolutions are minimal and secant, and no Koszul resolution in the box is secant, \Cref{thm:nosplit}, \Cref{thm:noci} \\
\texttt{transport\_}\allowbreak\texttt{scaling}, \texttt{transport\_}\allowbreak\texttt{multiplicity} & $\det\phi=c^{2G}$ and $\phi^{T}E\phi=c^{2}E$ on the samples; the multiplicity is $c^{4}$ on the stable sublattice and smaller on the unstable one, and the lower bound increases strictly, \Cref{lem:transport}, \Cref{lem:multiplicity} \\
\texttt{quaternionic\_}\allowbreak\texttt{ratios}, \texttt{quaternionic\_}\allowbreak\texttt{scaling}, \texttt{quaternionic\_}\allowbreak\texttt{presentation} & the scale-free ratios of the quaternionic Weil cycle, the growth of its data in $b$ alone, and the change of presentation by $\alpha$, \Cref{prop:quatdivisible} \\
\texttt{quaternionic\_}\allowbreak\texttt{weil\_}\allowbreak\texttt{lattice\_}\allowbreak\texttt{0}, \texttt{quaternionic\_}\allowbreak\texttt{weil\_}\allowbreak\texttt{lattice\_}\allowbreak\texttt{1}, \texttt{quaternionic\_}\allowbreak\texttt{weil\_}\allowbreak\texttt{lattice\_}\allowbreak\texttt{2}, \texttt{quaternionic\_}\allowbreak\texttt{weil\_}\allowbreak\texttt{lattice\_}\allowbreak\texttt{3} & for $d=1,2,3,7$ the integral Weil lattice has Gram matrix $\operatorname{diag}(8d,8d^{2})$ and the divisor products meet it in a sublattice of index $2(a_{1}a_{2})^{2}$, \Cref{prop:quatdivisible} \\
\texttt{divisor\_}\allowbreak\texttt{route\_}\allowbreak\texttt{rank}, \texttt{divisor\_}\allowbreak\texttt{route\_}\allowbreak\texttt{trace}, \texttt{divisor\_}\allowbreak\texttt{route\_}\allowbreak\texttt{pencil}, \texttt{divisor\_}\allowbreak\texttt{route\_}\allowbreak\texttt{index} & the lattice of $K$-bilinear classes has rank $12$ and signature $(8,4)$; the pencil $\phi_{k}^{2}=(1+k^{2})/(4d^{2})$ with $x_{k}$ primitive; and $\ZZ x_{k}+\ZZ x_{k}(i\cdot,\cdot)$ saturated for every $k$, \Cref{prop:pencil} \\
\texttt{divisor\_}\allowbreak\texttt{route\_}\allowbreak\texttt{minimum\_}\allowbreak\texttt{one}, \texttt{divisor\_}\allowbreak\texttt{route\_}\allowbreak\texttt{minimum\_}\allowbreak\texttt{three} & $N_{k}=\ZZ\tfrac{1}{2d}\eta\oplus\ZZ x_{k}\oplus\ZZ x_{k}(i\cdot,\cdot)$ and the minimum $\mu_{k}=32d^{2}(1+k^{2})$ of $I$ on it off $\QQ\eta$, for $d=1,3$ and $k\le10$, \Cref{prop:pencil} \\
\texttt{divisor\_}\allowbreak\texttt{route\_}\allowbreak\texttt{unitary\_}\allowbreak\texttt{one}, \texttt{divisor\_}\allowbreak\texttt{route\_}\allowbreak\texttt{unitary\_}\allowbreak\texttt{three}, \texttt{divisor\_}\allowbreak\texttt{route\_}\allowbreak\texttt{centraliser\_}\allowbreak\texttt{one\_}\allowbreak\texttt{0}, \texttt{divisor\_}\allowbreak\texttt{route\_}\allowbreak\texttt{centraliser\_}\allowbreak\texttt{one\_}\allowbreak\texttt{1}, \texttt{divisor\_}\allowbreak\texttt{route\_}\allowbreak\texttt{centraliser\_}\allowbreak\texttt{one\_}\allowbreak\texttt{2}, \texttt{divisor\_}\allowbreak\texttt{route\_}\allowbreak\texttt{centraliser\_}\allowbreak\texttt{one\_}\allowbreak\texttt{3}, \texttt{divisor\_}\allowbreak\texttt{route\_}\allowbreak\texttt{centraliser\_}\allowbreak\texttt{one\_}\allowbreak\texttt{4}, \texttt{divisor\_}\allowbreak\texttt{route\_}\allowbreak\texttt{centraliser\_}\allowbreak\texttt{three\_}\allowbreak\texttt{0}, \texttt{divisor\_}\allowbreak\texttt{route\_}\allowbreak\texttt{centraliser\_}\allowbreak\texttt{three\_}\allowbreak\texttt{1}, \texttt{divisor\_}\allowbreak\texttt{route\_}\allowbreak\texttt{centraliser\_}\allowbreak\texttt{three\_}\allowbreak\texttt{2}, \texttt{divisor\_}\allowbreak\texttt{route\_}\allowbreak\texttt{centraliser\_}\allowbreak\texttt{three\_}\allowbreak\texttt{3}, \texttt{divisor\_}\allowbreak\texttt{route\_}\allowbreak\texttt{centraliser\_}\allowbreak\texttt{three\_}\allowbreak\texttt{4} & the centraliser of $i$ in $\mathfrak{sp}(V,E)$ has dimension $16$, and that of $i$ and $\phi_{x_{k}}$ dimension $10$ with three invariants in $\wedge^{2}V^{*}$, for $d=1,3$ and $k\le4$, \Cref{prop:pencil} \\
\texttt{divisor\_}\allowbreak\texttt{route\_}\allowbreak\texttt{siegel\_}\allowbreak\texttt{one}, \texttt{divisor\_}\allowbreak\texttt{route\_}\allowbreak\texttt{siegel\_}\allowbreak\texttt{three} & a rational complex structure on the Siegel locus and its brackets span the centraliser, $d=1,3$, \Cref{prop:pencil} \\
\texttt{exceptional\_}\allowbreak\texttt{mumford\_}\allowbreak\texttt{square}, \texttt{exceptional\_}\allowbreak\texttt{quintic}, \texttt{exceptional\_}\allowbreak\texttt{annihilator} & two of the eight invariants of $\wedge^{4}(V\oplus V)$ on a Mumford square are exceptional; the adjoint weights of the quintic threefold; and the annihilator of the Weil class in $HT^{2}$ for $n=2,3$, \Cref{prop:mumfordproduct}, \Cref{prop:noabeliantype}, \Cref{thm:p2forces} \\
\texttt{cm\_}\allowbreak\texttt{fields\_}\allowbreak\texttt{quartic}, \texttt{cm\_}\allowbreak\texttt{fields\_}\allowbreak\texttt{sextic\_}\allowbreak\texttt{seven\_}\allowbreak\texttt{two}, \texttt{cm\_}\allowbreak\texttt{fields\_}\allowbreak\texttt{sextic\_}\allowbreak\texttt{seven\_}\allowbreak\texttt{one}, \texttt{cm\_}\allowbreak\texttt{fields\_}\allowbreak\texttt{sextic\_}\allowbreak\texttt{nine} & the base point for $\QQ(\zeta_{5})$, $\QQ(\zeta_{8})$, $\QQ(\zeta_{7})$ and $\QQ(\zeta_{9})$: CM type, trace-dual bases, eigenvectors, the alternating form and the closed form of the balanced classes, \Cref{thm:cmbasepoint} \\
\texttt{cm\_}\allowbreak\texttt{fields\_}\allowbreak\texttt{neron\_}\allowbreak\texttt{severi}, \texttt{cm\_}\allowbreak\texttt{fields\_}\allowbreak\texttt{composite} & $\NS$ has dimension $32$ at the base point for $\QQ(\zeta_{5})$, $n=2$; and $W_{K}$ lies in the span of the products of pairs of $W_{F}$ for $\QQ(i)\subset\QQ(\zeta_{8})$, \Cref{prop:composite} \\
\texttt{mumford\_}\allowbreak\texttt{rm\_}\allowbreak\texttt{isotypic}, \texttt{mumford\_}\allowbreak\texttt{rm\_}\allowbreak\texttt{isotypic\_}\allowbreak\texttt{mixed}, \texttt{mumford\_}\allowbreak\texttt{rm\_}\allowbreak\texttt{isotypic\_}\allowbreak\texttt{diagonal}, \texttt{mumford\_}\allowbreak\texttt{rm\_}\allowbreak\texttt{image} & $\wedge^{2}V$ has four summands of multiplicity one, the image of $\mu$ fills each, $h^{2,0}=6$, and the Hodge numbers of the pieces, \Cref{lem:mumfordmu}, \Cref{prop:mumfordrm} \\
\texttt{mumford\_}\allowbreak\texttt{rm\_}\allowbreak\texttt{pairing}, \texttt{mumford\_}\allowbreak\texttt{rm\_}\allowbreak\texttt{adjoint\_}\allowbreak\texttt{one}, \texttt{mumford\_}\allowbreak\texttt{rm\_}\allowbreak\texttt{adjoint\_}\allowbreak\texttt{two} & the Lefschetz pairing and the trace form agree up to scale, and $4\mu\mu^{\dagger}$ acts by scalars on the pieces, \Cref{prop:mumfordrm} \\
\texttt{mumford\_}\allowbreak\texttt{rm\_}\allowbreak\texttt{spin}, \texttt{mumford\_}\allowbreak\texttt{rm\_}\allowbreak\texttt{kuga\_}\allowbreak\texttt{satake} & the spin lifts, thirty-two vectors of highest weight $(1,1,1)$, and the Kuga--Satake map, with $N_{i}$ anticommuting and squaring to $4,8,20$, \Cref{prop:ksembed}, \Cref{thm:mumfordks} \\
\texttt{twistor\_}\allowbreak\texttt{real\_}\allowbreak\texttt{structure}, \texttt{twistor\_}\allowbreak\texttt{quaternion}, \texttt{twistor\_}\allowbreak\texttt{pieces}, \texttt{twistor\_}\allowbreak\texttt{exchange} & the real structure and the three complex structures of a Mumford square, the quaternion that rules out a polarisation at $J_{1}$, and the types and exchange of the tensors $\pi$, \Cref{prop:notwistor} \\
\texttt{twistor\_}\allowbreak\texttt{invariant\_}\allowbreak\texttt{one}, \texttt{twistor\_}\allowbreak\texttt{invariant\_}\allowbreak\texttt{two}, \texttt{twistor\_}\allowbreak\texttt{annihilator\_}\allowbreak\texttt{first}, \texttt{twistor\_}\allowbreak\texttt{annihilator\_}\allowbreak\texttt{third} & $\omega$ is invariant, and its annihilator at $J_{1}$ and $J_{3}$ is one line, \Cref{prop:notwistor} \\
\texttt{hk\_}\allowbreak\texttt{iota}, \texttt{hk\_}\allowbreak\texttt{isotypic}, \texttt{hk\_}\allowbreak\texttt{products}, \texttt{hk\_}\allowbreak\texttt{octic} & the map $\iota$ and the symplectic form, the isotypic pieces of $H^{2}(X\times X)$, the exceptional classes as products, and the octic form, \Cref{prop:hksquare} \\
\texttt{semireg\_}\allowbreak\texttt{monomials}, \texttt{semireg\_}\allowbreak\texttt{block} & the six products of $\beta,\widehat\beta,\ell$ are independent and $\eta^{2}$, $\theta^{+}\theta^{-}$ are independent; the multiplication table of one block in its Hodge basis, \Cref{lem:sixmonomials}, \Cref{lem:splitsemireg} \\
\texttt{semireg\_}\allowbreak\texttt{rank\_}\allowbreak\texttt{two}, \texttt{semireg\_}\allowbreak\texttt{rank\_}\allowbreak\texttt{three} & the semiregularity map of a sum of line bundles is injective for up to three classes at $n=2$ and ten at $n=3$, and has rank $22<24$ for four classes at $n=2$, \Cref{lem:tensorrank}, \Cref{thm:positivity} \\
\texttt{semireg\_}\allowbreak\texttt{object\_}\allowbreak\texttt{chern}, \texttt{semireg\_}\allowbreak\texttt{object\_}\allowbreak\texttt{rank} & the explicit object: its Chern character, index and Prouhet--Tarry--Escott pair, and the rank $75$ of its semiregularity map, \Cref{lem:pte}, \Cref{prop:explicitobject} \\
\texttt{qk\_}\allowbreak\texttt{model}, \texttt{qk\_}\allowbreak\texttt{closed\_}\allowbreak\texttt{form} & the one-place model: the transforms of $X\times0$, the diagonal, the antidiagonal and six graphs, and $\ch\Phi_{j}(T_{j}-P_{j})=-\frac14(\eta_{j}^{2}+\gamma_{j}^{2}-d\ell_{j}^{2})$ as a polynomial identity in $d$, a kernel not used in the paper \\
\texttt{qk\_}\allowbreak\texttt{su\_}\allowbreak\texttt{one}, \texttt{qk\_}\allowbreak\texttt{su\_}\allowbreak\texttt{two}, \texttt{qk\_}\allowbreak\texttt{su\_}\allowbreak\texttt{three}, \texttt{qk\_}\allowbreak\texttt{su\_}\allowbreak\texttt{five\_}\allowbreak\texttt{halves}, \texttt{qk\_}\allowbreak\texttt{su\_}\allowbreak\texttt{two\_}\allowbreak\texttt{fifths} & integral bases of $\mathfrak{su}_{j}(d)$ for $d=1,2,3,\frac52,\frac25$, \Cref{prop:quarticexclusions} \\
\texttt{qk\_}\allowbreak\texttt{flat}, \texttt{qk\_}\allowbreak\texttt{eigen\_}\allowbreak\texttt{hodge} & the class $\ch\Phi_{j}(T_{j}-P_{j})$ is flat, of pure degree four, with nonzero Weil part and of Hodge type, its twists by $e^{\pm\ell/2}$ are not flat, $\int\eta_{j}^{4}:\int\Omega_{j}^{2}=3:4$, and $\sqrt{-q}$ has eigenvalues $\pm4s_{j}$ on $(\gamma_{j}\mp s_{j}\ell_{j})^{2}$, not used in the paper \\
\texttt{qk\_}\allowbreak\texttt{contraction\_}\allowbreak\texttt{one}, \texttt{qk\_}\allowbreak\texttt{contraction\_}\allowbreak\texttt{one\_}\allowbreak\texttt{general}, \texttt{qk\_}\allowbreak\texttt{contraction\_}\allowbreak\texttt{three}, \texttt{qk\_}\allowbreak\texttt{contraction\_}\allowbreak\texttt{three\_}\allowbreak\texttt{general} & its contraction map has rank $23$ with a $5$-dimensional kernel in $H^{1}(T)$, against $24$ and $4$ for a general invariant class, at $d=1,3$, not used in the paper \\
\texttt{qk\_}\allowbreak\texttt{invariants\_}\allowbreak\texttt{one\_}\allowbreak\texttt{low}, \texttt{qk\_}\allowbreak\texttt{invariants\_}\allowbreak\texttt{one\_}\allowbreak\texttt{middle}, \texttt{qk\_}\allowbreak\texttt{invariants\_}\allowbreak\texttt{one\_}\allowbreak\texttt{high}, \texttt{qk\_}\allowbreak\texttt{commutant\_}\allowbreak\texttt{one}, \texttt{qk\_}\allowbreak\texttt{invariants\_}\allowbreak\texttt{three\_}\allowbreak\texttt{low}, \texttt{qk\_}\allowbreak\texttt{invariants\_}\allowbreak\texttt{three\_}\allowbreak\texttt{middle}, \texttt{qk\_}\allowbreak\texttt{invariants\_}\allowbreak\texttt{three\_}\allowbreak\texttt{high}, \texttt{qk\_}\allowbreak\texttt{commutant\_}\allowbreak\texttt{three}, \texttt{qk\_}\allowbreak\texttt{invariants\_}\allowbreak\texttt{two\_}\allowbreak\texttt{fifths\_}\allowbreak\texttt{low}, \texttt{qk\_}\allowbreak\texttt{invariants\_}\allowbreak\texttt{two\_}\allowbreak\texttt{fifths\_}\allowbreak\texttt{middle}, \texttt{qk\_}\allowbreak\texttt{invariants\_}\allowbreak\texttt{two\_}\allowbreak\texttt{fifths\_}\allowbreak\texttt{high}, \texttt{qk\_}\allowbreak\texttt{commutant\_}\allowbreak\texttt{two\_}\allowbreak\texttt{fifths} & for $d=1,3,\frac25$ the invariants of $\mathfrak{su}_{j}(d)$ have dimensions $1,1,3,1,1$ in degrees $0,2,4,6,8$ and none in odd degree, and the commutant is spanned by $1$ and $\sqrt{-q}$, \Cref{prop:quarticexclusions} \\
\texttt{qk\_}\allowbreak\texttt{graph\_}\allowbreak\texttt{classes\_}\allowbreak\texttt{one}, \texttt{qk\_}\allowbreak\texttt{graph\_}\allowbreak\texttt{classes\_}\allowbreak\texttt{two}, \texttt{qk\_}\allowbreak\texttt{graph\_}\allowbreak\texttt{classes\_}\allowbreak\texttt{three}, \texttt{qk\_}\allowbreak\texttt{graph\_}\allowbreak\texttt{classes\_}\allowbreak\texttt{four}, \texttt{qk\_}\allowbreak\texttt{graphs}, \texttt{qk\_}\allowbreak\texttt{graph\_}\allowbreak\texttt{invariants}, \texttt{qk\_}\allowbreak\texttt{graph\_}\allowbreak\texttt{twists} & the classes $(b,a)_{*}(c)$ on seven graphs span a space $L_{j}$ of dimension $24$, and $\Phi_{j}(L_{j})$ meets the invariants exactly in $\mathrm{span}(1,\mathrm{pt}_{j})$, and $e^{B}\Phi_{j}(L_{j})$ exactly in $e^{g\eta_{j}}\mathrm{span}(1,\mathrm{pt}_{j})$ for four $B$, \Cref{prop:quarticexclusions} \\
\texttt{qk\_}\allowbreak\texttt{pure\_}\allowbreak\texttt{spinors\_}\allowbreak\texttt{one}, \texttt{qk\_}\allowbreak\texttt{pure\_}\allowbreak\texttt{spinors\_}\allowbreak\texttt{two}, \texttt{qk\_}\allowbreak\texttt{pure\_}\allowbreak\texttt{spinors\_}\allowbreak\texttt{three}, \texttt{qk\_}\allowbreak\texttt{pure\_}\allowbreak\texttt{spinors\_}\allowbreak\texttt{four}, \texttt{qk\_}\allowbreak\texttt{pure\_}\allowbreak\texttt{spinors\_}\allowbreak\texttt{five}, \texttt{qk\_}\allowbreak\texttt{pure\_}\allowbreak\texttt{spinors\_}\allowbreak\texttt{six}, \texttt{qk\_}\allowbreak\texttt{pure\_}\allowbreak\texttt{spinors\_}\allowbreak\texttt{seven}, \texttt{qk\_}\allowbreak\texttt{not\_}\allowbreak\texttt{pure} & the transforms of line bundles on seven graphs are pure spinors, also after the twist by $e^{\ell/2}$; at $d=1$ the flat classes $\ch\Phi_{j}(T_{j}-P_{j})$ and $\Omega_{j}$ are not pure, and $(\gamma_{j}-i\ell_{j})^{2}$ is pure over $\CC$, \Cref{prop:quarticexclusions} \\
\texttt{qk\_}\allowbreak\texttt{parity\_}\allowbreak\texttt{congruences}, \texttt{qk\_}\allowbreak\texttt{mod\_}\allowbreak\texttt{small}, \texttt{qk\_}\allowbreak\texttt{congruences\_}\allowbreak\texttt{large}, \texttt{qk\_}\allowbreak\texttt{nu\_}\allowbreak\texttt{bounds} & $\int x^{2}$ is even for integral classes of degree four on a torus of dimension four; the conditions (a) to (c) leave only $D=3$, $m=1$ among all squarefree $D$; and the bounds on $\nu$, \Cref{rem:quarticnegext} \\
\texttt{qk\_}\allowbreak\texttt{prym}, \texttt{qk\_}\allowbreak\texttt{prym\_}\allowbreak\texttt{bound} & $\varphi(N)=4$ exactly for $N=5,8,10,12$ below $200$, a hyperbolic hermitian form of rank $2n$ has determinant $(-1)^{n}$, and $3g-3<\frac12\varphi(N)(g-1)^{2}$ for all $\varphi(N)\ge4$, $g\ge3$, \Cref{rem:quarticprym} \\
\end{longtable}
\endgroup

An element of $\ZZ[\sqrt{-d}]$ is represented by a pair of integers, with
multiplication and conjugation the evident operations. With twelve exceptions,
every statement in the table is a closed Boolean expression and is proved by
\texttt{decide}, which asks the kernel to evaluate it by reduction; for the
larger computations, those in exterior algebras among them, the form
\texttt{decide +kernel} hands the evaluation to the kernel directly. This is
the strongest check available and the least flexible: it works only over
finite ranges, which is why each of these theorems carries an explicit bound,
and why none of them can be the general statement that the paper proves.
Where the statement is a polynomial identity, the box is chosen large enough
in each variable for agreement on it to be the identity. The rank
certificates, those of \texttt{polarised\_}\allowbreak\texttt{number\_}\allowbreak\texttt{three}
and the four theorems after it and many of the theorems of Sections 56 to 81,
record exact vectors in the kernel and a family of images independent modulo
a prime $p$ at which the ring of definition has a homomorphism to
$\FF_{p}$: $p=998244353$ or $p=1000033$, at which $-1$ is a square, for
$\ZZ[\sqrt{-1}]$; $p=754974721$, with $\sqrt{-d}$ sent to a recorded root,
for $\ZZ[\sqrt{-d}]$; and $p=1000003$ for integer matrices. Reduction modulo
$p$ is a ring map, so the family is independent in characteristic zero, and
the two sizes add up to the dimension of the space.
The twelve exceptions,
\texttt{criterion\_}\allowbreak\texttt{endomorphisms\_n2},
\texttt{criterion\_}\allowbreak\texttt{endomorphisms\_n3},
\texttt{quartic\_}\allowbreak\texttt{rank\_}\allowbreak\texttt{two\_}\allowbreak\texttt{large\_}\allowbreak\texttt{odd\_}\allowbreak\texttt{part},
the second half of \texttt{orlov\_}\allowbreak\texttt{n4\_}\allowbreak\texttt{euler},
\texttt{divisor\_}\allowbreak\texttt{endo\_}\allowbreak\texttt{positive},
\texttt{divisor\_}\allowbreak\texttt{rank\_}\allowbreak\texttt{two\_}\allowbreak\texttt{mod3},
\texttt{lattice\_}\allowbreak\texttt{u3v\_}\allowbreak\texttt{criterion},
\texttt{burch\_}\allowbreak\texttt{rank\_}\allowbreak\texttt{two\_}\allowbreak\texttt{parity},
\texttt{qk\_}\allowbreak\texttt{mod\_}\allowbreak\texttt{small},
\texttt{qk\_}\allowbreak\texttt{congruences\_}\allowbreak\texttt{large},
\texttt{qk\_}\allowbreak\texttt{nu\_}\allowbreak\texttt{bounds} and
\texttt{qk\_}\allowbreak\texttt{prym\_}\allowbreak\texttt{bound},
are stated for all natural numbers and proved from the lemmas on natural
numbers in the core library, with \texttt{decide} used at most for a final
check over finitely many residues. In every case nothing is quoted: the
kernel recomputes each assertion from the definitions.

Two of the entries need a word on what they settle and what they do not.
\texttt{hodge\_dim\_count} verifies the bookkeeping of \Cref{thm:hdgcount}:
summing the contributions $(n,n)$, $(2n,0)$ and $(0,2n)$ over the degrees
gives three in the middle degree and one elsewhere. The representation theory
that identifies those contributions is in the proof of \Cref{thm:hdgcount} and
is not formalised. Likewise, \texttt{k3\_type\_only\_at\_one} verifies that
the equation $h^{2,0}=1$ has $n=1$ as its only solution once the three
possible values $0$, $n^{2}$ and $n(n-1)$ are known; the identification of
those values is \Cref{prop:normpart}. In both cases the division of labour is
the one followed throughout this appendix: the kernel settles the computation,
and the mathematics around it is written out and must be read.

\subsection{Failure of the obvious witness}

Two theorems of the file show that statements a careful reader would accept
are false. Both bear on \Cref{sec:characters}.

The first concerns the witness $\tau=1+\sqrt{-d}$ for \Cref{lem:charsdiffer}.
It is the obvious choice, and it does not work for every field:
\texttt{bad\_witness} shows that at $d=1$ the ratio $\tau/\bbar{\tau}$ is
$\sqrt{-1}$ and at $d=3$ it is a primitive cube root of unity, so for those
fields the two characters agree at this $\tau$ for suitable $n$. The lemma is
true, but not by that route; this is why its proof in \Cref{sec:characters}
uses only the finiteness of $\mu(K)$ and avoids witnesses altogether.

The second concerns the exceptional set of that witness at $d=3$. It is
$3\mid n$, and not $6\mid n$ as the order of $\mu(K)$ might suggest;
\texttt{bad\_witness\_exceptions} records the exact condition. Neither
statement affects \Cref{thm:character}. Both are slips of the kind that a
careful reading can miss and a kernel check catches, and the file keeps them
for that reason.

\subsection{Limits of the formalisation}

Everything else is used as mathematics and is not formalised here: the
cohomology of an abelian variety as an exterior algebra, the Hodge
decomposition, the Fourier--Mukai transform as a functor,
Grothendieck--Riemann--Roch, the Lefschetz hyperplane theorem, hard Lefschetz,
the invariant theory of $\SL_{2n}$, the Clifford algebra as a functor, and the
Lefschetz standard conjecture. Formalising them would require a library of
complex algebraic geometry that does not at present exist. What the Lean file
establishes is that the arithmetic of \Cref{tab:lean} is correct; where the
paper reasons, the reader must still read.

Some computations of \Cref{sec:verification} have no counterpart in the
file. The Macaulay2 computations of local $\Ext$ algebras and jet classes,
items (XXIX), (XXXI), (XL) and the local part of (LXVIII), the evaluation of
theta integrals in ball arithmetic in part (H) of item (LXXII), and the
search of item (XXII) over up to $4.29\times10^{14}$ sign patterns are beyond
a kernel with no libraries. The exact linear algebra of items (I), (XIII),
(XVI), (XVIII), (XIX), (XXIII), (XXIV), (XXXIV) to (XXXVII), (XLI) to
(XLIII), (XLVII) to (XLIX), (LI), (LIII), (LIV) and (LVII) to (LXIII) is in
Sections 56 to 81 of the file, with the rank certificates described above,
except for these parts, which stay with the exact computations: the
exhaustive searches of item (XIII); the other five shapes at $n=3$ and the
longer run at $n=4,5$ of item (LIII); parts (A) to (C) of item (LVIII),
identities in the exterior algebra on twenty-four generators; parts (B) to
(D) of item (LIX), which work over both places of explicit number fields;
part (F) of item (LX), the invariance of classes of $X\times X$ and
$X\times X\times B$ under $\mathfrak{sl}(V)$; and parts (D) and (F) to (H)
of item (LXI), identities at random rational points, invariants of
orthogonal groups, spans in CM fields of degree up to ten, and the
cohomology of generalised Kummer varieties sector by sector.
